% !TeX program = lualatex
% ALIUS Bulletin native visual reconstruction source.
% Generated from off-repo reference layout data; does not include or input a pre-existing article PDF.
\ifdefined\ALIUSIssueBuild
\else
\documentclass{article}
\usepackage[paperwidth=595.2756bp,paperheight=841.8898bp,margin=0bp]{geometry}
\usepackage{graphicx}
\usepackage{xcolor}
\usepackage{tikz}
\usepackage{iftex}
\usepackage[hidelinks]{hyperref}
\ifPDFTeX
  \usepackage[T1]{fontenc}
  \usepackage[utf8]{inputenc}
  \usepackage{textcomp}
  \usepackage{amssymb}
\else
  \usepackage{fontspec}
\fi
\pagestyle{empty}
\fi
\ifPDFTeX
  \providecommand{\ALIUSDeclareUnicodeFallbacks}{%
    \DeclareUnicodeCharacter{00AC}{\ensuremath{\neg}}%
    \DeclareUnicodeCharacter{00BE}{\ensuremath{\frac{3}{4}}}%
    \DeclareUnicodeCharacter{0101}{\=a}%
    \DeclareUnicodeCharacter{0107}{\'c}%
    \DeclareUnicodeCharacter{012B}{\=\i}%
    \DeclareUnicodeCharacter{0131}{\i}%
    \DeclareUnicodeCharacter{0142}{\l}%
    \DeclareUnicodeCharacter{015B}{\'s}%
    \DeclareUnicodeCharacter{016B}{\=u}%
    \DeclareUnicodeCharacter{025C}{q}%
    \DeclareUnicodeCharacter{0278}{t}%
    \DeclareUnicodeCharacter{02AE}{Th}%
    \DeclareUnicodeCharacter{02B0}{ff}%
    \DeclareUnicodeCharacter{02B4}{ffi}%
    \DeclareUnicodeCharacter{02BE}{ft}%
    \DeclareUnicodeCharacter{02BF}{fi}%
    \DeclareUnicodeCharacter{02C0}{fl}%
    \DeclareUnicodeCharacter{02C1}{fi}%
    \DeclareUnicodeCharacter{02D9}{Th}%
    \DeclareUnicodeCharacter{02EF}{gy}%
    \DeclareUnicodeCharacter{0394}{\ensuremath{\Delta}}%
    \DeclareUnicodeCharacter{03B2}{\ensuremath{\beta}}%
    \DeclareUnicodeCharacter{03BA}{\ensuremath{\kappa}}%
    \DeclareUnicodeCharacter{068D}{-}%
    \DeclareUnicodeCharacter{08BC}{ti}%
    \DeclareUnicodeCharacter{099E}{ti}%
    \DeclareUnicodeCharacter{1E43}{\d{m}}%
    \DeclareUnicodeCharacter{1E47}{\d{n}}%
    \DeclareUnicodeCharacter{1E63}{\d{s}}%
    \DeclareUnicodeCharacter{201C}{``}%
    \DeclareUnicodeCharacter{201D}{''}%
    \DeclareUnicodeCharacter{2260}{\ensuremath{\neq}}%
    \DeclareUnicodeCharacter{25A1}{\ensuremath{\square}}%
  }%
  \ifdefined\ALIUSIssueBuild\else\ALIUSDeclareUnicodeFallbacks\fi
  \providecommand{\ALIUSFontLatoLight}{\fontfamily{lato-LF}\fontseries{l}\selectfont}
  \providecommand{\ALIUSFontLatoRegular}{\fontfamily{lato-LF}\fontseries{m}\selectfont}
  \providecommand{\ALIUSFontLatoItalic}{\fontfamily{lato-LF}\fontseries{m}\itshape}
  \providecommand{\ALIUSFontLatoLightItalic}{\fontfamily{lato-LF}\fontseries{l}\itshape}
  \providecommand{\ALIUSFontCormorantRegular}{\fontfamily{CormorantGaramond-LF}\fontseries{m}\selectfont}
  \providecommand{\ALIUSFontCormorantLight}{\fontfamily{CormorantGaramond-LF}\fontseries{l}\selectfont}
  \providecommand{\ALIUSFontCormorantMedium}{\fontfamily{CormorantGaramond-LF}\fontseries{medium}\selectfont}
  \providecommand{\ALIUSFontCormorantBold}{\fontfamily{CormorantGaramond-LF}\fontseries{b}\selectfont}
  \providecommand{\ALIUSFontCormorantItalic}{\fontfamily{CormorantGaramond-LF}\fontseries{m}\itshape}
  \providecommand{\ALIUSFontTimes}{\fontfamily{ptm}\selectfont}
  \providecommand{\ALIUSFontCalibri}{\fontfamily{phv}\selectfont}
  \providecommand{\ALIUSFontCambria}{\fontfamily{ptm}\selectfont}
  \providecommand{\ALIUSFontSymbolFallback}{\fontfamily{ptm}\selectfont}
\else
  \providecommand{\ALIUSFontLatoLight}{\fontspec{Lato Light}}
  \providecommand{\ALIUSFontLatoRegular}{\fontspec{Lato Regular}}
  \providecommand{\ALIUSFontLatoItalic}{\fontspec{Lato Italic}}
  \providecommand{\ALIUSFontLatoLightItalic}{\fontspec{Lato Light Italic}}
  \providecommand{\ALIUSFontCormorantRegular}{\fontspec{Cormorant Garamond Regular}}
  \providecommand{\ALIUSFontCormorantLight}{\fontspec{Cormorant Garamond Light}}
  \providecommand{\ALIUSFontCormorantMedium}{\fontspec{Cormorant Garamond Medium}}
  \providecommand{\ALIUSFontCormorantBold}{\fontspec{Cormorant Garamond Bold}}
  \providecommand{\ALIUSFontCormorantItalic}{\fontspec{Cormorant Garamond Italic}}
  \providecommand{\ALIUSFontTimes}{\fontspec{Times New Roman}}
  \providecommand{\ALIUSFontCalibri}{\fontspec{Calibri}}
  \providecommand{\ALIUSFontCambria}{\fontspec{Cambria}}
  \providecommand{\ALIUSFontSymbolFallback}{\fontspec{Times New Roman}}
\fi
\providecommand{\ALIUSPullQuoteOpen}{\textquotedblleft}
\providecommand{\ALIUSPullQuoteClose}{\textquotedblright}
\definecolor{ALIUSC000000}{HTML}{000000}
\definecolor{ALIUSC0000FF}{HTML}{0000FF}
\definecolor{ALIUSC1F8135}{HTML}{1F8135}
\definecolor{ALIUSC595959}{HTML}{595959}
\definecolor{ALIUSC767171}{HTML}{767171}
\definecolor{ALIUSCBFBFBF}{HTML}{BFBFBF}
\providecommand{\ALIUSPlacedText}[6]{%
  \node[anchor=north west,inner sep=0pt,outer sep=0pt,text=#4] at (#1bp,-#2bp) {%
    \resizebox{#3bp}{!}{{#5\fontsize{#6bp}{#6bp}\selectfont #6\kern0pt}}%
  };%
}
\providecommand{\ALIUSPlacedTextContent}[7]{%
  \node[anchor=north west,inner sep=0pt,outer sep=0pt,text=#4] at (#1bp,-#2bp) {%
    \resizebox{#3bp}{!}{{#5\fontsize{#6bp}{#6bp}\selectfont #7}}%
  };%
}
\ifdefined\ALIUSIssueBuild
\else
\begin{document}
\fi
% Page 1
\thispagestyle{empty}
\noindent\begin{tikzpicture}[remember picture,overlay,x=1bp,y=1bp,shift={(current page.north west)}]
  \fill[ALIUSC1F8135] (372.711bp,-159.591bp) rectangle ++(78.480bp,-0.240bp);
  \fill[ALIUSC1F8135] (372.711bp,-219.591bp) rectangle ++(98.160bp,-0.240bp);
  \fill[ALIUSC1F8135] (372.711bp,-275.991bp) rectangle ++(96.720bp,-0.240bp);
  \fill[ALIUSCBFBFBF] (367.191bp,-136.551bp) rectangle ++(0.240bp,-88.800bp);
  \fill[ALIUSC0000FF] (88.791bp,-286.311bp) rectangle ++(165.840bp,-0.480bp);
  \fill[ALIUSCBFBFBF] (367.191bp,-225.351bp) rectangle ++(0.240bp,-2.160bp);
  \fill[ALIUSCBFBFBF] (367.191bp,-227.511bp) rectangle ++(0.240bp,-71.040bp);
  \ALIUSPlacedTextContent{90.951}{72.970}{339.004}{ALIUSC000000}{\ALIUSFontLatoRegular}{19.920}{The Neuroscience of Decision Making}
  \ALIUSPlacedTextContent{372.711}{136.454}{76.167}{ALIUSC000000}{\ALIUSFontLatoRegular}{11.040}{Timothy Hanks}
  \ALIUSPlacedTextContent{372.711}{149.629}{83.371}{ALIUSC1F8135}{\ALIUSFontLatoLight}{9.120}{thanks@ucdavis.edu}
  \ALIUSPlacedTextContent{88.791}{155.720}{130.764}{ALIUSC000000}{\ALIUSFontLatoLight}{16.080}{An interview with}
  \ALIUSPlacedTextContent{372.711}{160.429}{100.578}{ALIUSC000000}{\ALIUSFontLatoLight}{9.120}{Center for Neuroscience,}
  \ALIUSPlacedTextContent{372.711}{171.229}{77.157}{ALIUSC000000}{\ALIUSFontLatoLight}{9.120}{UC Davis, CA, USA}
  \ALIUSPlacedTextContent{88.791}{174.957}{129.036}{ALIUSC000000}{\ALIUSFontLatoLight}{18.960}{Timothy Hanks}
  \ALIUSPlacedTextContent{372.711}{196.454}{106.867}{ALIUSC000000}{\ALIUSFontLatoRegular}{11.040}{Alexandra Mikhailova}
  \ALIUSPlacedTextContent{372.711}{209.629}{100.556}{ALIUSC1F8135}{\ALIUSFontLatoLight}{9.120}{amikhailova@ucdavis.edu}
  \ALIUSPlacedTextContent{88.791}{209.679}{263.481}{ALIUSC000000}{\ALIUSFontLatoLight}{12.960}{by Alexandra Mikhailova \& Daniel A. Friedman}
  \ALIUSPlacedTextContent{372.711}{220.429}{100.578}{ALIUSC000000}{\ALIUSFontLatoLight}{9.120}{Center for Neuroscience,}
  % ALIUS normalized citation panel begin
  % ALIUS normalized citation panel text: Cite as: Hanks, T., Mikhailova, A., \& Friedman, D.A. (2020). The Neuroscience of Decision Making. An interview with Timothy Hanks by Alexandra Mikhailova and Daniel A. Friedman. ALIUS Bulletin, 4, 68-80,
  \fill[white] (87.291bp,-225.402bp) rectangle ++(276.920bp,-95.500bp);
  \node[anchor=north west,inner sep=0pt,outer sep=0pt,text=ALIUSC000000] at (88.791bp,-227.402bp) {\begin{minipage}{273.920bp}{\ALIUSFontLatoLight\fontsize{9.200bp}{10.700bp}\selectfont\raggedright Cite as: Hanks, T., Mikhailova, A., \& Friedman, D.A. (2020). The Neuroscience of Decision Making. An interview with Timothy Hanks by Alexandra Mikhailova and Daniel A. Friedman. ALIUS Bulletin, 4, 68-80,\par}\end{minipage}};
  \node[anchor=south east,inner sep=0pt,outer sep=0pt,text=ALIUSC1F8135] at (362.711bp,-316.402bp) {{\ALIUSFontLatoLight\fontsize{8.600bp}{8.600bp}\selectfont\href{https://doi.org/10.34700/8pg4-0h12}{https://doi.org/10.34700/8pg4-0h12}}};
  % ALIUS normalized citation panel end
  \ALIUSPlacedTextContent{372.711}{231.229}{77.157}{ALIUSC000000}{\ALIUSFontLatoLight}{9.120}{UC Davis, CA, USA}
  \ALIUSPlacedTextContent{372.711}{252.854}{94.999}{ALIUSC000000}{\ALIUSFontLatoRegular}{11.040}{Daniel A. Friedman}
  \ALIUSPlacedTextContent{372.711}{266.029}{99.247}{ALIUSC1F8135}{\ALIUSFontLatoLight}{9.120}{dafriedman@ucdavis.edu}
  \ALIUSPlacedTextContent{372.711}{276.829}{111.643}{ALIUSC000000}{\ALIUSFontLatoLight}{9.120}{Department of Entomology}
  \ALIUSPlacedTextContent{372.711}{287.629}{137.613}{ALIUSC000000}{\ALIUSFontLatoLight}{9.120}{\& Nematology, UC Davis, CA, USA}
  \ALIUSPlacedTextContent{90.951}{328.732}{55.719}{ALIUSC000000}{\ALIUSFontLatoRegular}{13.920}{Abstract}
  \ALIUSPlacedTextContent{90.951}{357.267}{418.272}{ALIUSC000000}{\ALIUSFontLatoLight}{12.000}{In this interview, Professor Tim Hanks discusses topics related to neuroscience,}
  \ALIUSPlacedTextContent{90.951}{371.667}{418.272}{ALIUSC000000}{\ALIUSFontLatoLight}{12.000}{decision making, philosophy, and science as a career. Hanks explores how ideas}
  \ALIUSPlacedTextContent{90.951}{386.067}{418.278}{ALIUSC000000}{\ALIUSFontLatoLight}{12.000}{from computational neuroscience have helped him set his own research agenda}
  \ALIUSPlacedTextContent{90.951}{400.467}{418.270}{ALIUSC000000}{\ALIUSFontLatoLight}{12.000}{and also navigate everyday situations. The way the brain makes decisions is}
  \ALIUSPlacedTextContent{90.951}{414.867}{418.270}{ALIUSC000000}{\ALIUSFontLatoLight}{12.000}{deeply intertwined with topics such as free will, conscious awareness, and mental}
  \ALIUSPlacedTextContent{90.951}{429.267}{418.268}{ALIUSC000000}{\ALIUSFontLatoLight}{12.000}{health. In order to productively study such diverse topics related to decision}
  \ALIUSPlacedTextContent{90.951}{443.667}{418.268}{ALIUSC000000}{\ALIUSFontLatoLight}{12.000}{making, Hanks recommends an integrative approach that draws on multiple types}
  \ALIUSPlacedTextContent{90.951}{458.067}{418.271}{ALIUSC000000}{\ALIUSFontLatoLight}{12.000}{of experiments and model systems, with an eye towards clinical deployment. His}
  \ALIUSPlacedTextContent{90.951}{472.467}{418.270}{ALIUSC000000}{\ALIUSFontLatoLight}{12.000}{approach builds on various scientific frameworks, while also reminding us to stay}
  \ALIUSPlacedTextContent{90.951}{486.867}{395.636}{ALIUSC000000}{\ALIUSFontLatoLight}{12.000}{open-minded about what the future of neuroscience may look like or bring.}
  \ALIUSPlacedTextContent{90.951}{513.267}{47.340}{ALIUSC000000}{\ALIUSFontLatoItalic}{12.000}{keywords}
  \ALIUSPlacedTextContent{138.289}{513.267}{302.416}{ALIUSC000000}{\ALIUSFontLatoLightItalic}{12.000}{: decision making, Bayesian, free will, mental health, attention}
  \ALIUSPlacedTextContent{90.951}{554.373}{418.377}{ALIUSC1F8135}{\ALIUSFontLatoLight}{12.357}{To quote your book chapter from 'Neurobiology of Decision Making': "Thus}
  \ALIUSPlacedTextContent{90.951}{571.653}{418.354}{ALIUSC1F8135}{\ALIUSFontLatoLight}{12.357}{looking for the right questions is just another kind of decision" (Shadlen et al.,}
  \ALIUSPlacedTextContent{90.951}{588.933}{418.362}{ALIUSC1F8135}{\ALIUSFontLatoLight}{12.357}{2008). This quote resonates with folk sayings such as "the answer is in the}
  \ALIUSPlacedTextContent{90.951}{606.213}{418.299}{ALIUSC1F8135}{\ALIUSFontLatoLight}{12.357}{question", and also with computational perspectives on decision making that}
  \ALIUSPlacedTextContent{90.951}{623.253}{418.400}{ALIUSC1F8135}{\ALIUSFontLatoLight}{12.357}{highlight the role of (beliefs about) prior beliefs in controlling how evidential}
  \ALIUSPlacedTextContent{90.951}{640.533}{418.307}{ALIUSC1F8135}{\ALIUSFontLatoLight}{12.357}{stimuli play a role in decision making. As a scientist, how did you come to study}
  \ALIUSPlacedTextContent{90.951}{657.813}{418.337}{ALIUSC1F8135}{\ALIUSFontLatoLight}{12.357}{decision making? What decision making rules or heuristics were helpful for}
  \ALIUSPlacedTextContent{90.951}{675.093}{341.628}{ALIUSC1F8135}{\ALIUSFontLatoLight}{12.357}{you during your development as a professional and researcher?}
  \ALIUSPlacedTextContent{90.951}{704.249}{418.266}{ALIUSC000000}{\ALIUSFontCormorantRegular}{13.920}{I was always drawn to the questions that I found most mysterious, with those}
  \ALIUSPlacedTextContent{90.951}{723.688}{418.327}{ALIUSC000000}{\ALIUSFontCormorantRegular}{13.920}{of mental experience among the top of the list. These are the questions that}
  \ALIUSPlacedTextContent{90.951}{743.128}{421.533}{ALIUSC000000}{\ALIUSFontCormorantRegular}{13.920}{I find most engaging, that keep me thinking at night. What}
  \ALIUSPlacedTextContent{90.951}{793.143}{122.397}{ALIUSC767171}{\ALIUSFontCormorantLight}{12.000}{ALIUS Bulletin n°4 (2020)}
  \ALIUSPlacedTextContent{387.435}{793.143}{120.728}{ALIUSC767171}{\ALIUSFontCormorantLight}{12.000}{aliusresearch.org/bulletin}
\end{tikzpicture}
\null
\newpage
% Page 2
\thispagestyle{empty}
\noindent\begin{tikzpicture}[remember picture,overlay,x=1bp,y=1bp,shift={(current page.north west)}]
  \ALIUSPlacedTextContent{90.951}{36.329}{321.512}{ALIUSC767171}{\ALIUSFontCormorantRegular}{13.920}{Timothy Hanks – The Neuroscience of Decision Making}
  \ALIUSPlacedTextContent{492.951}{36.329}{16.277}{ALIUSC000000}{\ALIUSFontCormorantRegular}{13.920}{69}
  \ALIUSPlacedTextContent{107.791}{83.799}{14.160}{ALIUSC595959}{\ALIUSFontCormorantLight}{48.000}{\ALIUSPullQuoteOpen}
  \ALIUSPlacedTextContent{135.151}{91.867}{322.809}{ALIUSC000000}{\ALIUSFontLatoLight}{15.120}{The most helpful decision-making advice I have}
  \ALIUSPlacedTextContent{135.151}{109.867}{322.791}{ALIUSC000000}{\ALIUSFontLatoLight}{15.120}{applied in my life and career is to consider every}
  \ALIUSPlacedTextContent{466.591}{122.919}{14.160}{ALIUSC595959}{\ALIUSFontCormorantLight}{48.000}{\ALIUSPullQuoteClose}
  \ALIUSPlacedTextContent{135.151}{127.867}{171.083}{ALIUSC000000}{\ALIUSFontLatoLight}{15.120}{choice in a Bayesian sense}
  \ALIUSPlacedTextContent{90.951}{179.609}{418.274}{ALIUSC000000}{\ALIUSFontCormorantRegular}{13.920}{provided the final push in the direction of neuroscience was the belief that}
  \ALIUSPlacedTextContent{90.951}{199.049}{418.270}{ALIUSC000000}{\ALIUSFontCormorantRegular}{13.920}{these questions were ripe for finding answers. For decision making in}
  \ALIUSPlacedTextContent{90.951}{218.488}{418.190}{ALIUSC000000}{\ALIUSFontCormorantRegular}{13.920}{particular, I was drawn in by a fascination with the specific research being}
  \ALIUSPlacedTextContent{90.951}{238.169}{418.291}{ALIUSC000000}{\ALIUSFontCormorantRegular}{13.920}{conducted in that area that I learned about as an undergraduate student.}
  \ALIUSPlacedTextContent{90.951}{257.608}{418.235}{ALIUSC000000}{\ALIUSFontCormorantRegular}{13.920}{That's what led me to do pre- and post-doctoral research in the labs of Mike}
  \ALIUSPlacedTextContent{90.951}{277.048}{418.196}{ALIUSC000000}{\ALIUSFontCormorantRegular}{13.920}{Shadlen and Carlos Brody. I was enthusiastic enough about both of their}
  \ALIUSPlacedTextContent{90.951}{296.488}{418.286}{ALIUSC000000}{\ALIUSFontCormorantRegular}{13.920}{research programs that I can still recall in vivid detail the joy of positive email}
  \ALIUSPlacedTextContent{90.951}{316.168}{418.268}{ALIUSC000000}{\ALIUSFontCormorantRegular}{13.920}{correspondence with both before joining their labs. So, I followed the path}
  \ALIUSPlacedTextContent{90.951}{335.608}{418.318}{ALIUSC000000}{\ALIUSFontCormorantRegular}{13.920}{of my strongest interests. The most important advice I can give is to follow a}
  \ALIUSPlacedTextContent{90.951}{355.048}{418.214}{ALIUSC000000}{\ALIUSFontCormorantRegular}{13.920}{path that results in sustained day-to-day internal satisfaction, and to take}
  \ALIUSPlacedTextContent{90.951}{374.488}{327.350}{ALIUSC000000}{\ALIUSFontCormorantRegular}{13.920}{efforts to be honest with yourself about what that involves.}
  \ALIUSPlacedTextContent{90.951}{406.168}{418.185}{ALIUSC000000}{\ALIUSFontCormorantRegular}{13.920}{Along my path, I have also carefully considered what would be most}
  \ALIUSPlacedTextContent{90.951}{425.608}{418.233}{ALIUSC000000}{\ALIUSFontCormorantRegular}{13.920}{beneficial in the long term to allow me to continue to follow my interests.}
  \ALIUSPlacedTextContent{90.951}{445.048}{418.254}{ALIUSC000000}{\ALIUSFontCormorantRegular}{13.920}{That's a lot harder to determine. I think the most helpful decision-making}
  \ALIUSPlacedTextContent{90.951}{464.488}{418.218}{ALIUSC000000}{\ALIUSFontCormorantRegular}{13.920}{advice I have applied in my life and career is to consider every choice in a}
  \ALIUSPlacedTextContent{90.951}{484.168}{418.160}{ALIUSC000000}{\ALIUSFontCormorantRegular}{13.920}{Bayesian sense by trying to use evidence to estimate probabilities or at the}
  \ALIUSPlacedTextContent{90.951}{503.608}{418.143}{ALIUSC000000}{\ALIUSFontCormorantRegular}{13.920}{very least, to explicitly consider uncertainty. This can obviously help to avoid}
  \ALIUSPlacedTextContent{90.951}{523.049}{418.273}{ALIUSC000000}{\ALIUSFontCormorantRegular}{13.920}{overconfidence, but it can also help to overcome paralyzing doubt. Most}
  \ALIUSPlacedTextContent{90.951}{542.489}{418.215}{ALIUSC000000}{\ALIUSFontCormorantRegular}{13.920}{importantly, it gives a principled foundation for choosing courses of action.}
  \ALIUSPlacedTextContent{90.951}{561.928}{354.342}{ALIUSC000000}{\ALIUSFontCormorantRegular}{13.920}{And that's ultimately why we have a brain, to shape our actions.}
  \ALIUSPlacedTextContent{90.951}{593.733}{418.343}{ALIUSC1F8135}{\ALIUSFontLatoLight}{12.357}{Every day we are consciously aware of deliberative decisions that we make}
  \ALIUSPlacedTextContent{90.951}{611.013}{418.226}{ALIUSC1F8135}{\ALIUSFontLatoLight}{12.357}{(e.g. which clothes to wear, what to write in an email to a colleague), rising like}
  \ALIUSPlacedTextContent{90.951}{628.053}{418.392}{ALIUSC1F8135}{\ALIUSFontLatoLight}{12.357}{islands of awareness from a vast ocean of subconscious decisions (e.g.}
  \ALIUSPlacedTextContent{90.951}{645.333}{418.388}{ALIUSC1F8135}{\ALIUSFontLatoLight}{12.357}{physiological decisions related to blood pressure, heart rate, oculomotor}
  \ALIUSPlacedTextContent{90.951}{662.613}{414.986}{ALIUSC1F8135}{\ALIUSFontLatoLight}{12.357}{tracking, postural balance). There are also situations where subconsciously-}
  \ALIUSPlacedTextContent{90.951}{679.893}{418.362}{ALIUSC1F8135}{\ALIUSFontLatoLight}{12.357}{decided preferences can apparently be steered by conscious input, for}
  \ALIUSPlacedTextContent{90.951}{697.173}{418.452}{ALIUSC1F8135}{\ALIUSFontLatoLight}{12.357}{example in the case of a deliberative pacific response to an aggressive}
  \ALIUSPlacedTextContent{90.951}{714.453}{418.375}{ALIUSC1F8135}{\ALIUSFontLatoLight}{12.357}{stimulus. How do you think about the role of conscious awareness in decision}
  \ALIUSPlacedTextContent{90.951}{731.733}{418.432}{ALIUSC1F8135}{\ALIUSFontLatoLight}{12.357}{making across these domains, and especially in the case of decisions where}
  \ALIUSPlacedTextContent{90.951}{749.013}{324.032}{ALIUSC1F8135}{\ALIUSFontLatoLight}{12.357}{phenomenologically we seem to "have a say" in the output?}
  \ALIUSPlacedTextContent{90.951}{793.143}{122.397}{ALIUSC767171}{\ALIUSFontCormorantLight}{12.000}{ALIUS Bulletin n°4 (2020)}
  \ALIUSPlacedTextContent{387.435}{793.143}{120.728}{ALIUSC767171}{\ALIUSFontCormorantLight}{12.000}{aliusresearch.org/bulletin}
\end{tikzpicture}
\null
\newpage
% Page 3
\thispagestyle{empty}
\noindent\begin{tikzpicture}[remember picture,overlay,x=1bp,y=1bp,shift={(current page.north west)}]
  \ALIUSPlacedTextContent{90.951}{36.329}{321.512}{ALIUSC767171}{\ALIUSFontCormorantRegular}{13.920}{Timothy Hanks – The Neuroscience of Decision Making}
  \ALIUSPlacedTextContent{493.191}{36.329}{15.941}{ALIUSC000000}{\ALIUSFontCormorantRegular}{13.920}{70}
  \ALIUSPlacedTextContent{90.951}{73.173}{418.310}{ALIUSC1F8135}{\ALIUSFontLatoLight}{12.357}{What neural events might distinguish the awareness of a decision from ones}
  \ALIUSPlacedTextContent{90.951}{90.453}{418.369}{ALIUSC1F8135}{\ALIUSFontLatoLight}{12.357}{made subconsciously? Is conscious awareness necessary, sufficient, or}
  \ALIUSPlacedTextContent{90.951}{107.733}{418.393}{ALIUSC1F8135}{\ALIUSFontLatoLight}{12.357}{differentially associated with decision making tasks (Ganupuru et al., 2019) or}
  \ALIUSPlacedTextContent{90.951}{125.013}{418.482}{ALIUSC1F8135}{\ALIUSFontLatoLight}{12.357}{meta-cognitive assessments? How do we coherently generate research}
  \ALIUSPlacedTextContent{90.951}{142.293}{418.452}{ALIUSC1F8135}{\ALIUSFontLatoLight}{12.357}{agendas that try to bridge the experience of conscious awareness, with neural}
  \ALIUSPlacedTextContent{90.951}{159.333}{158.248}{ALIUSC1F8135}{\ALIUSFontLatoLight}{12.357}{or molecular measurements?}
  \ALIUSPlacedTextContent{90.951}{188.488}{418.249}{ALIUSC000000}{\ALIUSFontCormorantRegular}{13.920}{Let me start with the "have a say" component of the question because that is}
  \ALIUSPlacedTextContent{90.951}{207.928}{418.115}{ALIUSC000000}{\ALIUSFontCormorantRegular}{13.920}{where I think we have the best opportunity for scientific leverage. This relates}
  \ALIUSPlacedTextContent{90.951}{227.609}{418.200}{ALIUSC000000}{\ALIUSFontCormorantRegular}{13.920}{to the question of free will, which I think is often posed in a way that stifles}
  \ALIUSPlacedTextContent{90.951}{247.049}{418.272}{ALIUSC000000}{\ALIUSFontCormorantRegular}{13.920}{progress on finding answers. If the question is, "Do we have free will?", the}
  \ALIUSPlacedTextContent{90.951}{266.488}{418.154}{ALIUSC000000}{\ALIUSFontCormorantRegular}{13.920}{answer depends on how we define free will. I think a better starting point is}
  \ALIUSPlacedTextContent{90.951}{285.928}{418.261}{ALIUSC000000}{\ALIUSFontCormorantRegular}{13.920}{to ask the question, "To what extent do we have free will?" This promotes}
  \ALIUSPlacedTextContent{90.951}{305.608}{418.159}{ALIUSC000000}{\ALIUSFontCormorantRegular}{13.920}{inquiry of what it is that we actually have. As a scientist studying decision}
  \ALIUSPlacedTextContent{90.951}{325.048}{418.235}{ALIUSC000000}{\ALIUSFontCormorantRegular}{13.920}{making, I can see a clear roadmap towards answering this by describing the}
  \ALIUSPlacedTextContent{90.951}{344.488}{418.296}{ALIUSC000000}{\ALIUSFontCormorantRegular}{13.920}{neural mechanisms that underlie our decisions. It's my belief that when that}
  \ALIUSPlacedTextContent{90.951}{363.928}{418.243}{ALIUSC000000}{\ALIUSFontCormorantRegular}{13.920}{description is complete, it will provide a satisfying and compelling answer to}
  \ALIUSPlacedTextContent{90.951}{383.608}{418.289}{ALIUSC000000}{\ALIUSFontCormorantRegular}{13.920}{the question of free will much in the way that our current understanding of}
  \ALIUSPlacedTextContent{90.951}{403.048}{418.245}{ALIUSC000000}{\ALIUSFontCormorantRegular}{13.920}{biology, including genetics and evolutionary theory, provide a satisfying and}
  \ALIUSPlacedTextContent{90.951}{422.488}{228.161}{ALIUSC000000}{\ALIUSFontCormorantRegular}{13.920}{compelling answer to the question of life.}
  \ALIUSPlacedTextContent{90.951}{453.928}{418.296}{ALIUSC000000}{\ALIUSFontCormorantRegular}{13.920}{To make this more concrete, let me give some examples of what I think these}
  \ALIUSPlacedTextContent{90.951}{473.608}{418.345}{ALIUSC000000}{\ALIUSFontCormorantRegular}{13.920}{types of explanations could look like. Through the work of many researchers}
  \ALIUSPlacedTextContent{90.951}{493.048}{418.151}{ALIUSC000000}{\ALIUSFontCormorantRegular}{13.920}{across our field, we have found that commitment to a decision in many}
  \ALIUSPlacedTextContent{90.951}{512.489}{418.273}{ALIUSC000000}{\ALIUSFontCormorantRegular}{13.920}{situations can be described as occurring through the accumulation of}
  \ALIUSPlacedTextContent{90.951}{531.928}{418.256}{ALIUSC000000}{\ALIUSFontCormorantRegular}{13.920}{evidence to a threshold level or bound (Gold \& Shadlen, 2007; Hanks \&}
  \ALIUSPlacedTextContent{90.951}{551.609}{418.235}{ALIUSC000000}{\ALIUSFontCormorantRegular}{13.920}{Summerfield, 2017). The evidence is represented in neural activity and the}
  \ALIUSPlacedTextContent{90.951}{571.048}{418.257}{ALIUSC000000}{\ALIUSFontCormorantRegular}{13.920}{bound has been hypothesized to be applied to the level of neural activity. A}
  \ALIUSPlacedTextContent{90.951}{590.489}{418.296}{ALIUSC000000}{\ALIUSFontCormorantRegular}{13.920}{higher bound would require more evidence for decision commitment and a}
  \ALIUSPlacedTextContent{90.951}{609.928}{418.192}{ALIUSC000000}{\ALIUSFontCormorantRegular}{13.920}{lower bound would require less evidence. In other words, a person could}
  \ALIUSPlacedTextContent{90.951}{629.609}{418.257}{ALIUSC000000}{\ALIUSFontCormorantRegular}{13.920}{respond differently to the same evidence depending on the level of the}
  \ALIUSPlacedTextContent{117.151}{662.679}{14.160}{ALIUSC595959}{\ALIUSFontCormorantLight}{48.000}{\ALIUSPullQuoteOpen}
  \ALIUSPlacedTextContent{149.071}{673.147}{301.717}{ALIUSC000000}{\ALIUSFontLatoLight}{15.120}{Perhaps in following this roadmap, we will find}
  \ALIUSPlacedTextContent{149.071}{691.147}{52.583}{ALIUSC000000}{\ALIUSFontLatoLight}{15.120}{aspects}
  \ALIUSPlacedTextContent{214.099}{691.147}{17.542}{ALIUSC000000}{\ALIUSFontLatoLight}{15.120}{of}
  \ALIUSPlacedTextContent{244.086}{691.147}{56.909}{ALIUSC000000}{\ALIUSFontLatoLight}{15.120}{decision}
  \ALIUSPlacedTextContent{313.440}{691.147}{49.950}{ALIUSC000000}{\ALIUSFontLatoLight}{15.120}{making}
  \ALIUSPlacedTextContent{375.594}{691.147}{30.034}{ALIUSC000000}{\ALIUSFontLatoLight}{15.120}{that}
  \ALIUSPlacedTextContent{417.832}{691.147}{32.907}{ALIUSC000000}{\ALIUSFontLatoLight}{15.120}{defy}
  \ALIUSPlacedTextContent{149.071}{709.147}{301.650}{ALIUSC000000}{\ALIUSFontLatoLight}{15.120}{explanation based on our current scientific}
  \ALIUSPlacedTextContent{149.071}{727.147}{301.777}{ALIUSC000000}{\ALIUSFontLatoLight}{15.120}{understanding, or that are better explained}
  \ALIUSPlacedTextContent{466.831}{738.039}{14.160}{ALIUSC595959}{\ALIUSFontCormorantLight}{48.000}{\ALIUSPullQuoteClose}
  \ALIUSPlacedTextContent{149.071}{745.147}{297.737}{ALIUSC000000}{\ALIUSFontLatoLight}{15.120}{through new scientific conceptual frameworks}
  \ALIUSPlacedTextContent{90.951}{793.143}{122.397}{ALIUSC767171}{\ALIUSFontCormorantLight}{12.000}{ALIUS Bulletin n°4 (2020)}
  \ALIUSPlacedTextContent{387.435}{793.143}{120.728}{ALIUSC767171}{\ALIUSFontCormorantLight}{12.000}{aliusresearch.org/bulletin}
\end{tikzpicture}
\null
\newpage
% Page 4
\thispagestyle{empty}
\noindent\begin{tikzpicture}[remember picture,overlay,x=1bp,y=1bp,shift={(current page.north west)}]
  \ALIUSPlacedTextContent{90.951}{36.329}{321.512}{ALIUSC767171}{\ALIUSFontCormorantRegular}{13.920}{Timothy Hanks – The Neuroscience of Decision Making}
  \ALIUSPlacedTextContent{495.351}{36.329}{13.911}{ALIUSC000000}{\ALIUSFontCormorantRegular}{13.920}{71}
  \ALIUSPlacedTextContent{90.951}{73.048}{418.247}{ALIUSC000000}{\ALIUSFontCormorantRegular}{13.920}{decision bound. The bound is one mechanism for exerting will.}
  \ALIUSPlacedTextContent{90.951}{92.488}{418.282}{ALIUSC000000}{\ALIUSFontCormorantRegular}{13.920}{Understanding the factors that go into how a person sets their decision}
  \ALIUSPlacedTextContent{90.951}{111.928}{421.533}{ALIUSC000000}{\ALIUSFontCormorantRegular}{13.920}{bound then tells us something about how that person's will can be set.}
  \ALIUSPlacedTextContent{90.951}{131.609}{418.210}{ALIUSC000000}{\ALIUSFontCormorantRegular}{13.920}{This is just one example. There are many other neural mechanisms for}
  \ALIUSPlacedTextContent{90.951}{151.049}{418.245}{ALIUSC000000}{\ALIUSFontCormorantRegular}{13.920}{exerting will in decision making. I think that when we have discovered and}
  \ALIUSPlacedTextContent{90.951}{170.488}{418.257}{ALIUSC000000}{\ALIUSFontCormorantRegular}{13.920}{understood those mechanisms, that will provide a compelling account of the}
  \ALIUSPlacedTextContent{90.951}{189.928}{185.295}{ALIUSC000000}{\ALIUSFontCormorantRegular}{13.920}{extent to which we have free will.}
  \ALIUSPlacedTextContent{90.951}{221.609}{418.423}{ALIUSC000000}{\ALIUSFontCormorantRegular}{13.920}{In reading this, one might think I am describing an approach aimed at}
  \ALIUSPlacedTextContent{90.951}{241.049}{418.181}{ALIUSC000000}{\ALIUSFontCormorantRegular}{13.920}{yielding a deterministic account that will be reducible to our current}
  \ALIUSPlacedTextContent{90.951}{260.488}{418.221}{ALIUSC000000}{\ALIUSFontCormorantRegular}{13.920}{understanding of physical laws, but I don't think this is guaranteed. Perhaps}
  \ALIUSPlacedTextContent{90.951}{279.928}{418.207}{ALIUSC000000}{\ALIUSFontCormorantRegular}{13.920}{in following this roadmap, we will find aspects of decision making that defy}
  \ALIUSPlacedTextContent{90.951}{299.608}{418.237}{ALIUSC000000}{\ALIUSFontCormorantRegular}{13.920}{explanation based on our current scientific understanding, or that are better}
  \ALIUSPlacedTextContent{90.951}{319.048}{314.316}{ALIUSC000000}{\ALIUSFontCormorantRegular}{13.920}{explained through new scientific conceptual frameworks.}
  \ALIUSPlacedTextContent{90.951}{350.488}{418.204}{ALIUSC000000}{\ALIUSFontCormorantRegular}{13.920}{This is where I suspect consciousness will come to the fore. Currently, we}
  \ALIUSPlacedTextContent{90.951}{369.928}{418.263}{ALIUSC000000}{\ALIUSFontCormorantRegular}{13.920}{have no satisfying and compelling scientific explanation of consciousness.}
  \ALIUSPlacedTextContent{90.951}{389.608}{418.269}{ALIUSC000000}{\ALIUSFontCormorantRegular}{13.920}{Yet, as your question suggests, our internal experience gives us the impression}
  \ALIUSPlacedTextContent{90.951}{409.048}{418.330}{ALIUSC000000}{\ALIUSFontCormorantRegular}{13.920}{that consciousness matters for our decisions. If that impression is correct,}
  \ALIUSPlacedTextContent{90.951}{428.488}{418.339}{ALIUSC000000}{\ALIUSFontCormorantRegular}{13.920}{then in trying to fully explain the mechanisms of our decisions, I suspect that}
  \ALIUSPlacedTextContent{90.951}{447.928}{418.325}{ALIUSC000000}{\ALIUSFontCormorantRegular}{13.920}{eventually we will hit that wall. The hope is that how we hit the wall will}
  \ALIUSPlacedTextContent{90.951}{467.608}{418.316}{ALIUSC000000}{\ALIUSFontCormorantRegular}{13.920}{reveal how our current scientific conceptions need to be refined to better}
  \ALIUSPlacedTextContent{90.951}{487.048}{145.900}{ALIUSC000000}{\ALIUSFontCormorantRegular}{13.920}{understand consciousness.}
  \ALIUSPlacedTextContent{90.951}{518.489}{418.324}{ALIUSC000000}{\ALIUSFontCormorantRegular}{13.920}{As we pursue this scientific path, I think it is important to keep in mind that}
  \ALIUSPlacedTextContent{90.951}{537.928}{418.151}{ALIUSC000000}{\ALIUSFontCormorantRegular}{13.920}{consciousness is not a singular type. In this respect, I find it useful to consider}
  \ALIUSPlacedTextContent{90.951}{557.369}{418.289}{ALIUSC000000}{\ALIUSFontCormorantRegular}{13.920}{the evolutionary history of consciousness. Are different types of}
  \ALIUSPlacedTextContent{90.951}{577.048}{418.215}{ALIUSC000000}{\ALIUSFontCormorantRegular}{13.920}{consciousness (e.g., visual, auditory, volitional) the product of divergent or}
  \ALIUSPlacedTextContent{90.951}{596.489}{418.290}{ALIUSC000000}{\ALIUSFontCormorantRegular}{13.920}{convergent evolution? It seems possible that different forms of consciousness}
  \ALIUSPlacedTextContent{90.951}{615.928}{418.317}{ALIUSC000000}{\ALIUSFontCormorantRegular}{13.920}{derived from convergent evolution with distinct types of relationships with}
  \ALIUSPlacedTextContent{90.951}{635.369}{418.143}{ALIUSC000000}{\ALIUSFontCormorantRegular}{13.920}{neural or molecular measurements. Even if they derived from something}
  \ALIUSPlacedTextContent{90.951}{655.048}{418.213}{ALIUSC000000}{\ALIUSFontCormorantRegular}{13.920}{more akin to divergent evolution from a common proto-consciousness, I}
  \ALIUSPlacedTextContent{90.951}{674.489}{418.257}{ALIUSC000000}{\ALIUSFontCormorantRegular}{13.920}{would still not be surprised if there were distinct types of relationships with}
  \ALIUSPlacedTextContent{90.951}{693.928}{418.268}{ALIUSC000000}{\ALIUSFontCormorantRegular}{13.920}{neural or molecular measurements for different forms of consciousness. That}
  \ALIUSPlacedTextContent{90.951}{713.369}{418.290}{ALIUSC000000}{\ALIUSFontCormorantRegular}{13.920}{could prove very helpful in our quest for deeper understanding because the}
  \ALIUSPlacedTextContent{90.951}{733.049}{421.533}{ALIUSC000000}{\ALIUSFontCormorantRegular}{13.920}{commonalities between these relationships would potentially reveal}
  \ALIUSPlacedTextContent{90.951}{752.489}{193.203}{ALIUSC000000}{\ALIUSFontCormorantRegular}{13.920}{principles of a more general nature.}
  \ALIUSPlacedTextContent{90.951}{793.143}{122.397}{ALIUSC767171}{\ALIUSFontCormorantLight}{12.000}{ALIUS Bulletin n°4 (2020)}
  \ALIUSPlacedTextContent{387.435}{793.143}{120.728}{ALIUSC767171}{\ALIUSFontCormorantLight}{12.000}{aliusresearch.org/bulletin}
\end{tikzpicture}
\null
\newpage
% Page 5
\thispagestyle{empty}
\noindent\begin{tikzpicture}[remember picture,overlay,x=1bp,y=1bp,shift={(current page.north west)}]
  \ALIUSPlacedTextContent{90.951}{36.329}{321.512}{ALIUSC767171}{\ALIUSFontCormorantRegular}{13.920}{Timothy Hanks – The Neuroscience of Decision Making}
  \ALIUSPlacedTextContent{494.391}{36.329}{14.891}{ALIUSC000000}{\ALIUSFontCormorantRegular}{13.920}{72}
  \ALIUSPlacedTextContent{109.711}{77.319}{14.160}{ALIUSC595959}{\ALIUSFontCormorantLight}{48.000}{\ALIUSPullQuoteOpen}
  \ALIUSPlacedTextContent{137.311}{87.547}{322.773}{ALIUSC000000}{\ALIUSFontLatoLight}{15.120}{As we pursue this scientific path, I think it is}
  \ALIUSPlacedTextContent{137.311}{105.547}{322.727}{ALIUSC000000}{\ALIUSFontLatoLight}{15.120}{important to keep in mind that consciousness is}
  \ALIUSPlacedTextContent{470.911}{113.319}{14.160}{ALIUSC595959}{\ALIUSFontCormorantLight}{48.000}{\ALIUSPullQuoteClose}
  \ALIUSPlacedTextContent{137.311}{123.547}{119.989}{ALIUSC000000}{\ALIUSFontLatoLight}{15.120}{not a singular type}
  \ALIUSPlacedTextContent{90.951}{175.413}{418.363}{ALIUSC1F8135}{\ALIUSFontLatoLight}{12.357}{Your work on the speed vs. accuracy trade-off in the brains of macaques was}
  \ALIUSPlacedTextContent{90.951}{192.693}{418.461}{ALIUSC1F8135}{\ALIUSFontLatoLight}{12.357}{testing the hypothesis that fast decisions require a lower threshold of}
  \ALIUSPlacedTextContent{90.951}{209.973}{418.356}{ALIUSC1F8135}{\ALIUSFontLatoLight}{12.357}{evidence (Hanks et al., 2014). Your findings suggested that control of speed}
  \ALIUSPlacedTextContent{90.951}{227.013}{418.368}{ALIUSC1F8135}{\ALIUSFontLatoLight}{12.357}{vs. accuracy may be exerted through changes in decision related neural}
  \ALIUSPlacedTextContent{90.951}{244.293}{418.206}{ALIUSC1F8135}{\ALIUSFontLatoLight}{12.357}{activity itself, rather than through changes in the threshold applied to such}
  \ALIUSPlacedTextContent{90.951}{261.573}{418.345}{ALIUSC1F8135}{\ALIUSFontLatoLight}{12.357}{neural activity to terminate a decision. This is consistent with Bayesian}
  \ALIUSPlacedTextContent{90.951}{278.853}{418.401}{ALIUSC1F8135}{\ALIUSFontLatoLight}{12.357}{frameworks of decision making in the brain - i.e. that variation in internal}
  \ALIUSPlacedTextContent{90.951}{296.133}{418.277}{ALIUSC1F8135}{\ALIUSFontLatoLight}{12.357}{neural states or priors can weigh on the outcome of decisions, given the same}
  \ALIUSPlacedTextContent{90.951}{313.413}{418.436}{ALIUSC1F8135}{\ALIUSFontLatoLight}{12.357}{task stimuli. How do you think about the statistical or computational}
  \ALIUSPlacedTextContent{90.951}{330.693}{418.335}{ALIUSC1F8135}{\ALIUSFontLatoLight}{12.357}{perspectives on decision making, in regard to decision processes that can be}
  \ALIUSPlacedTextContent{90.951}{347.973}{418.386}{ALIUSC1F8135}{\ALIUSFontLatoLight}{12.357}{measured in the lab, as well as in the context of the complex long-term}
  \ALIUSPlacedTextContent{90.951}{365.013}{200.973}{ALIUSC1F8135}{\ALIUSFontLatoLight}{12.357}{decisions that humans have to make?}
  \ALIUSPlacedTextContent{90.951}{394.168}{418.122}{ALIUSC000000}{\ALIUSFontCormorantRegular}{13.920}{Yes, that is very nicely put, and connects directly to the point I was making}
  \ALIUSPlacedTextContent{90.951}{413.848}{418.307}{ALIUSC000000}{\ALIUSFontCormorantRegular}{13.920}{above. We studied the neural mechanisms that govern the tradeoff between}
  \ALIUSPlacedTextContent{90.951}{433.288}{418.184}{ALIUSC000000}{\ALIUSFontCormorantRegular}{13.920}{the speed and accuracy of decision making because we hope that they provide}
  \ALIUSPlacedTextContent{90.951}{452.728}{418.113}{ALIUSC000000}{\ALIUSFontCormorantRegular}{13.920}{more general insight into flexible control of decision commitment. In theory,}
  \ALIUSPlacedTextContent{90.951}{472.168}{418.230}{ALIUSC000000}{\ALIUSFontCormorantRegular}{13.920}{control over the speed vs. accuracy trade-off could occur in a variety of ways.}
  \ALIUSPlacedTextContent{90.951}{491.848}{418.229}{ALIUSC000000}{\ALIUSFontCormorantRegular}{13.920}{A lower "bound" for decision commitment might be implemented directly}
  \ALIUSPlacedTextContent{90.951}{511.288}{418.324}{ALIUSC000000}{\ALIUSFontCormorantRegular}{13.920}{through less neural activity being needed to trigger a choice. Alternatively, it}
  \ALIUSPlacedTextContent{90.951}{530.729}{418.057}{ALIUSC000000}{\ALIUSFontCormorantRegular}{13.920}{might be implemented with the same level of neural activity needed to trigger}
  \ALIUSPlacedTextContent{90.951}{550.169}{418.237}{ALIUSC000000}{\ALIUSFontCormorantRegular}{13.920}{a choice, but additional internally-generated drive to this neural activity. In}
  \ALIUSPlacedTextContent{90.951}{569.609}{418.212}{ALIUSC000000}{\ALIUSFontCormorantRegular}{13.920}{the both cases, there is less external drive needed, and therefore decisions are}
  \ALIUSPlacedTextContent{90.951}{589.289}{418.215}{ALIUSC000000}{\ALIUSFontCormorantRegular}{13.920}{made more hastily. We found the mechanism to be more internal drive, what}
  \ALIUSPlacedTextContent{90.951}{608.729}{418.346}{ALIUSC000000}{\ALIUSFontCormorantRegular}{13.920}{is sometimes referred to as an "urgency" signal to respond separate from the}
  \ALIUSPlacedTextContent{90.951}{628.169}{53.208}{ALIUSC000000}{\ALIUSFontCormorantRegular}{13.920}{evidence.}
  \ALIUSPlacedTextContent{90.951}{659.609}{418.174}{ALIUSC000000}{\ALIUSFontCormorantRegular}{13.920}{One thing that I really like about this result is that it seems to be a fairly}
  \ALIUSPlacedTextContent{90.951}{679.289}{418.190}{ALIUSC000000}{\ALIUSFontCormorantRegular}{13.920}{general mechanism for cognitive control of the brain. For decision making,}
  \ALIUSPlacedTextContent{90.951}{698.729}{418.271}{ALIUSC000000}{\ALIUSFontCormorantRegular}{13.920}{it has been replicated in multiple other studies from different labs and}
  \ALIUSPlacedTextContent{90.951}{718.169}{418.250}{ALIUSC000000}{\ALIUSFontCormorantRegular}{13.920}{different decision tasks (Heitz et al., 2014; Thura \& Cisek, 2016). There are}
  \ALIUSPlacedTextContent{90.951}{737.609}{421.533}{ALIUSC000000}{\ALIUSFontCormorantRegular}{13.920}{other situations where neuroscientists have described boosted neural activity}
  \ALIUSPlacedTextContent{90.951}{793.143}{122.397}{ALIUSC767171}{\ALIUSFontCormorantLight}{12.000}{ALIUS Bulletin n°4 (2020)}
  \ALIUSPlacedTextContent{387.435}{793.143}{120.728}{ALIUSC767171}{\ALIUSFontCormorantLight}{12.000}{aliusresearch.org/bulletin}
\end{tikzpicture}
\null
\newpage
% Page 6
\thispagestyle{empty}
\noindent\begin{tikzpicture}[remember picture,overlay,x=1bp,y=1bp,shift={(current page.north west)}]
  \ALIUSPlacedTextContent{90.951}{36.329}{321.512}{ALIUSC767171}{\ALIUSFontCormorantRegular}{13.920}{Timothy Hanks – The Neuroscience of Decision Making}
  \ALIUSPlacedTextContent{494.391}{36.329}{14.709}{ALIUSC000000}{\ALIUSFontCormorantRegular}{13.920}{73}
  \ALIUSPlacedTextContent{117.151}{77.319}{14.160}{ALIUSC595959}{\ALIUSFontCormorantLight}{48.000}{\ALIUSPullQuoteOpen}
  \ALIUSPlacedTextContent{145.471}{88.747}{306.659}{ALIUSC000000}{\ALIUSFontLatoLight}{15.120}{Even if you only care about the computational}
  \ALIUSPlacedTextContent{145.471}{106.747}{306.464}{ALIUSC000000}{\ALIUSFontLatoLight}{15.120}{type of understanding or only care about the}
  \ALIUSPlacedTextContent{145.471}{124.747}{104.643}{ALIUSC000000}{\ALIUSFontLatoLight}{15.120}{implementation}
  \ALIUSPlacedTextContent{267.368}{124.747}{32.913}{ALIUSC000000}{\ALIUSFontLatoLight}{15.120}{type}
  \ALIUSPlacedTextContent{317.536}{124.747}{17.545}{ALIUSC000000}{\ALIUSFontLatoLight}{15.120}{of}
  \ALIUSPlacedTextContent{352.094}{124.747}{99.888}{ALIUSC000000}{\ALIUSFontLatoLight}{15.120}{understanding,}
  \ALIUSPlacedTextContent{145.471}{142.747}{306.567}{ALIUSC000000}{\ALIUSFontLatoLight}{15.120}{studying both will probably help you more}
  \ALIUSPlacedTextContent{462.271}{155.559}{14.160}{ALIUSC595959}{\ALIUSFontCormorantLight}{48.000}{\ALIUSPullQuoteClose}
  \ALIUSPlacedTextContent{145.471}{160.747}{225.730}{ALIUSC000000}{\ALIUSFontLatoLight}{15.120}{effectively understand either alone}
  \ALIUSPlacedTextContent{90.951}{209.848}{418.257}{ALIUSC000000}{\ALIUSFontCormorantRegular}{13.920}{associated with cognitive control, perhaps the most prominent being with}
  \ALIUSPlacedTextContent{90.951}{229.288}{418.341}{ALIUSC000000}{\ALIUSFontCormorantRegular}{13.920}{selective attention (Squire et al., 2013). In that case, attention to a particular}
  \ALIUSPlacedTextContent{90.951}{248.729}{418.260}{ALIUSC000000}{\ALIUSFontCormorantRegular}{13.920}{part or feature of one's environment can boost activity of neurons selective}
  \ALIUSPlacedTextContent{90.951}{268.168}{418.196}{ALIUSC000000}{\ALIUSFontCormorantRegular}{13.920}{for that part or feature. While we don't know exactly how this activity is}
  \ALIUSPlacedTextContent{90.951}{287.848}{418.257}{ALIUSC000000}{\ALIUSFontCormorantRegular}{13.920}{boosted for attention or decision making, it's intriguing to consider that they}
  \ALIUSPlacedTextContent{90.951}{307.288}{418.115}{ALIUSC000000}{\ALIUSFontCormorantRegular}{13.920}{may be related mechanisms. Studying these mechanisms will help us know}
  \ALIUSPlacedTextContent{90.951}{326.728}{418.062}{ALIUSC000000}{\ALIUSFontCormorantRegular}{13.920}{more about what can be controlled through cognition and how. For example,}
  \ALIUSPlacedTextContent{90.951}{346.168}{418.080}{ALIUSC000000}{\ALIUSFontCormorantRegular}{13.920}{let's say that after it is started, the boosting can't be suppressed for some}
  \ALIUSPlacedTextContent{90.951}{365.848}{418.357}{ALIUSC000000}{\ALIUSFontCormorantRegular}{13.920}{period of time due to a limitation of the mechanism. In that case, we could}
  \ALIUSPlacedTextContent{90.951}{385.288}{418.357}{ALIUSC000000}{\ALIUSFontCormorantRegular}{13.920}{truly say that a person does not have the free will to reduce their decision}
  \ALIUSPlacedTextContent{90.951}{404.728}{418.215}{ALIUSC000000}{\ALIUSFontCormorantRegular}{13.920}{bound and increase it again within that time limiting span. I pose that merely}
  \ALIUSPlacedTextContent{90.951}{424.168}{418.206}{ALIUSC000000}{\ALIUSFontCormorantRegular}{13.920}{as an example. I do not expect things to be so clear cut. Instead, we will}
  \ALIUSPlacedTextContent{90.951}{443.848}{418.314}{ALIUSC000000}{\ALIUSFontCormorantRegular}{13.920}{probably find greater flexibility in cognitive control than this example, but}
  \ALIUSPlacedTextContent{90.951}{463.288}{397.615}{ALIUSC000000}{\ALIUSFontCormorantRegular}{13.920}{it may be the case that some forms of it are easier to achieve than others.}
  \ALIUSPlacedTextContent{90.951}{494.728}{418.185}{ALIUSC000000}{\ALIUSFontCormorantRegular}{13.920}{To address the final part of your question, the Bayesian framework provides}
  \ALIUSPlacedTextContent{90.951}{514.169}{418.215}{ALIUSC000000}{\ALIUSFontCormorantRegular}{13.920}{a theoretical foundation for understanding decision making, and we relate}
  \ALIUSPlacedTextContent{90.951}{533.849}{418.305}{ALIUSC000000}{\ALIUSFontCormorantRegular}{13.920}{much of our work to it (Beck et al., 2008). This "computational" perspective}
  \ALIUSPlacedTextContent{90.951}{553.289}{418.236}{ALIUSC000000}{\ALIUSFontCormorantRegular}{13.920}{is often contrasted with the brain's "implementation" within neural circuits.}
  \ALIUSPlacedTextContent{90.951}{572.729}{418.393}{ALIUSC000000}{\ALIUSFontCormorantRegular}{13.920}{When considering this distinction, it is important to realize that the}
  \ALIUSPlacedTextContent{90.951}{592.169}{418.252}{ALIUSC000000}{\ALIUSFontCormorantRegular}{13.920}{boundary may not always be clear cut and that knowledge of either can be}
  \ALIUSPlacedTextContent{90.951}{611.849}{418.261}{ALIUSC000000}{\ALIUSFontCormorantRegular}{13.920}{helpful for understanding both. In other words, even if you only care about}
  \ALIUSPlacedTextContent{90.951}{631.289}{418.307}{ALIUSC000000}{\ALIUSFontCormorantRegular}{13.920}{the computational type of understanding or only care about the}
  \ALIUSPlacedTextContent{90.951}{650.729}{418.190}{ALIUSC000000}{\ALIUSFontCormorantRegular}{13.920}{implementation type of understanding, studying both will probably help you}
  \ALIUSPlacedTextContent{90.951}{670.169}{226.118}{ALIUSC000000}{\ALIUSFontCormorantRegular}{13.920}{more effectively understand either alone.}
  \ALIUSPlacedTextContent{90.951}{701.973}{418.403}{ALIUSC1F8135}{\ALIUSFontLatoLight}{12.357}{Your lab developed a new change detection task based on a stream of}
  \ALIUSPlacedTextContent{90.951}{719.013}{418.374}{ALIUSC1F8135}{\ALIUSFontLatoLight}{12.357}{auditory clicks generated by a stochastic Poisson process (Johnson et al.,}
  \ALIUSPlacedTextContent{90.951}{736.293}{418.504}{ALIUSC1F8135}{\ALIUSFontLatoLight}{12.357}{2017), which allows you to measure the temporal weighting of sensory}
  \ALIUSPlacedTextContent{90.951}{753.573}{418.388}{ALIUSC1F8135}{\ALIUSFontLatoLight}{12.357}{evidence when subjects employ different decision-making rules (thus}
  \ALIUSPlacedTextContent{90.951}{793.143}{122.397}{ALIUSC767171}{\ALIUSFontCormorantLight}{12.000}{ALIUS Bulletin n°4 (2020)}
  \ALIUSPlacedTextContent{387.435}{793.143}{120.728}{ALIUSC767171}{\ALIUSFontCormorantLight}{12.000}{aliusresearch.org/bulletin}
\end{tikzpicture}
\null
\newpage
% Page 7
\thispagestyle{empty}
\noindent\begin{tikzpicture}[remember picture,overlay,x=1bp,y=1bp,shift={(current page.north west)}]
  \ALIUSPlacedTextContent{90.951}{36.329}{321.512}{ALIUSC767171}{\ALIUSFontCormorantRegular}{13.920}{Timothy Hanks – The Neuroscience of Decision Making}
  \ALIUSPlacedTextContent{493.671}{36.329}{15.577}{ALIUSC000000}{\ALIUSFontCormorantRegular}{13.920}{74}
  \ALIUSPlacedTextContent{90.951}{73.173}{418.389}{ALIUSC1F8135}{\ALIUSFontLatoLight}{12.357}{finessing the trade-off between false-positives and false-negatives). The}
  \ALIUSPlacedTextContent{90.951}{90.453}{418.327}{ALIUSC1F8135}{\ALIUSFontLatoLight}{12.357}{article concluded that "changes in decision stopping rules did not alter the}
  \ALIUSPlacedTextContent{90.951}{107.733}{418.448}{ALIUSC1F8135}{\ALIUSFontLatoLight}{12.357}{temporal weighting of sensory evidence on the decision in a systematic way.}
  \ALIUSPlacedTextContent{90.951}{125.013}{418.384}{ALIUSC1F8135}{\ALIUSFontLatoLight}{12.357}{Instead, it altered the magnitude of evidence needed to trigger a choice."}
  \ALIUSPlacedTextContent{90.951}{142.293}{418.534}{ALIUSC1F8135}{\ALIUSFontLatoLight}{12.357}{Unlike a forced-choice task, where participants have a window of time to}
  \ALIUSPlacedTextContent{90.951}{159.333}{418.427}{ALIUSC1F8135}{\ALIUSFontLatoLight}{12.357}{make a selection, your task allows subjects to "decide when to decide". What}
  \ALIUSPlacedTextContent{90.951}{176.613}{418.392}{ALIUSC1F8135}{\ALIUSFontLatoLight}{12.357}{other applications of this task would like to explore? What do you think is}
  \ALIUSPlacedTextContent{90.951}{193.893}{418.452}{ALIUSC1F8135}{\ALIUSFontLatoLight}{12.357}{gained or lost when we isolate decision making outcomes (e.g. bets) from the}
  \ALIUSPlacedTextContent{90.951}{211.173}{418.386}{ALIUSC1F8135}{\ALIUSFontLatoLight}{12.357}{time required to make such decisions? How do you think the methodological}
  \ALIUSPlacedTextContent{90.951}{228.453}{418.375}{ALIUSC1F8135}{\ALIUSFontLatoLight}{12.357}{tension between fixed-time and free-time decision making tasks could be}
  \ALIUSPlacedTextContent{90.951}{245.733}{234.122}{ALIUSC1F8135}{\ALIUSFontLatoLight}{12.357}{retrospectively or prospectively integrated?}
  \ALIUSPlacedTextContent{90.951}{274.888}{418.363}{ALIUSC000000}{\ALIUSFontCormorantRegular}{13.920}{We began to study change detection tasks because it allowed us to address an}
  \ALIUSPlacedTextContent{90.951}{294.328}{418.257}{ALIUSC000000}{\ALIUSFontCormorantRegular}{13.920}{important set of questions neglected by much of the previous work in the}
  \ALIUSPlacedTextContent{90.951}{313.768}{418.329}{ALIUSC000000}{\ALIUSFontCormorantRegular}{13.920}{field, including my own previous work. For many of the most common tasks}
  \ALIUSPlacedTextContent{90.951}{333.208}{418.163}{ALIUSC000000}{\ALIUSFontCormorantRegular}{13.920}{used to study the neural mechanisms of decision making, evidence should be}
  \ALIUSPlacedTextContent{90.951}{352.888}{418.218}{ALIUSC000000}{\ALIUSFontCormorantRegular}{13.920}{weighted equally across time. One of the major questions has been to}
  \ALIUSPlacedTextContent{90.951}{372.328}{418.204}{ALIUSC000000}{\ALIUSFontCormorantRegular}{13.920}{understand how that consistent weighting of new evidence can be achieved,}
  \ALIUSPlacedTextContent{90.951}{391.768}{418.243}{ALIUSC000000}{\ALIUSFontCormorantRegular}{13.920}{and tremendous progress has been made with answering that (Brody \&}
  \ALIUSPlacedTextContent{90.951}{411.208}{418.329}{ALIUSC000000}{\ALIUSFontCormorantRegular}{13.920}{Hanks, 2016). However, there are many situations where consistent weighting}
  \ALIUSPlacedTextContent{90.951}{430.888}{418.264}{ALIUSC000000}{\ALIUSFontCormorantRegular}{13.920}{of evidence across time is detrimental. In any situation with change or}
  \ALIUSPlacedTextContent{90.951}{450.328}{418.332}{ALIUSC000000}{\ALIUSFontCormorantRegular}{13.920}{instability, recent evidence is more informative than evidence from the more}
  \ALIUSPlacedTextContent{90.951}{469.768}{418.251}{ALIUSC000000}{\ALIUSFontCormorantRegular}{13.920}{distant past. So, there can actually be an advantage to "forgetting" – or at}
  \ALIUSPlacedTextContent{90.951}{489.208}{418.249}{ALIUSC000000}{\ALIUSFontCormorantRegular}{13.920}{least, not relying as much on information gathered further in the past (Glaze}
  \ALIUSPlacedTextContent{90.951}{508.648}{258.247}{ALIUSC000000}{\ALIUSFontCormorantRegular}{13.920}{et al., 2015; Radillo et al., 2017; Piet et al., 2018).}
  \ALIUSPlacedTextContent{90.951}{540.328}{418.289}{ALIUSC000000}{\ALIUSFontCormorantRegular}{13.920}{In trying to take the next step to understand how our brain weighs evidence}
  \ALIUSPlacedTextContent{90.951}{559.768}{418.403}{ALIUSC000000}{\ALIUSFontCormorantRegular}{13.920}{across time, we have been systematically exploring and characterizing our}
  \ALIUSPlacedTextContent{90.951}{579.208}{418.304}{ALIUSC000000}{\ALIUSFontCormorantRegular}{13.920}{capacity to evaluate evidence on different timescales. Clearly, people are}
  \ALIUSPlacedTextContent{90.951}{598.648}{418.293}{ALIUSC000000}{\ALIUSFontCormorantRegular}{13.920}{capable of changing the timescale of evidence evaluation for decision making.}
  \ALIUSPlacedTextContent{90.951}{618.328}{418.199}{ALIUSC000000}{\ALIUSFontCormorantRegular}{13.920}{But what are the limits of this flexibility? This is what we are trying to}
  \ALIUSPlacedTextContent{90.951}{637.768}{418.312}{ALIUSC000000}{\ALIUSFontCormorantRegular}{13.920}{determine. We have been using tasks that push this flexibility to see if people}
  \ALIUSPlacedTextContent{90.951}{657.208}{418.132}{ALIUSC000000}{\ALIUSFontCormorantRegular}{13.920}{can simultaneously evaluate the same sources of evidence across different}
  \ALIUSPlacedTextContent{90.951}{676.648}{418.247}{ALIUSC000000}{\ALIUSFontCormorantRegular}{13.920}{timescales. If you are making a quick decision about one aspect of something,}
  \ALIUSPlacedTextContent{90.951}{696.328}{418.385}{ALIUSC000000}{\ALIUSFontCormorantRegular}{13.920}{does it limit your ability to combine information simultaneously over longer}
  \ALIUSPlacedTextContent{90.951}{715.768}{418.137}{ALIUSC000000}{\ALIUSFontCormorantRegular}{13.920}{periods of time to make a more careful decision about some other aspect of}
  \ALIUSPlacedTextContent{90.951}{735.208}{87.844}{ALIUSC000000}{\ALIUSFontCormorantRegular}{13.920}{the same thing?}
  \ALIUSPlacedTextContent{90.951}{793.143}{122.397}{ALIUSC767171}{\ALIUSFontCormorantLight}{12.000}{ALIUS Bulletin n°4 (2020)}
  \ALIUSPlacedTextContent{387.435}{793.143}{120.728}{ALIUSC767171}{\ALIUSFontCormorantLight}{12.000}{aliusresearch.org/bulletin}
\end{tikzpicture}
\null
\newpage
% Page 8
\thispagestyle{empty}
\noindent\begin{tikzpicture}[remember picture,overlay,x=1bp,y=1bp,shift={(current page.north west)}]
  \ALIUSPlacedTextContent{90.951}{36.329}{321.512}{ALIUSC767171}{\ALIUSFontCormorantRegular}{13.920}{Timothy Hanks – The Neuroscience of Decision Making}
  \ALIUSPlacedTextContent{494.151}{36.329}{14.989}{ALIUSC000000}{\ALIUSFontCormorantRegular}{13.920}{75}
  \ALIUSPlacedTextContent{90.951}{73.048}{418.257}{ALIUSC000000}{\ALIUSFontCormorantRegular}{13.920}{These questions seem most important in situations where we must determine}
  \ALIUSPlacedTextContent{90.951}{92.488}{414.974}{ALIUSC000000}{\ALIUSFontCormorantRegular}{13.920}{when to decide, but I would argue that many "forced-choice" decision-}
  \ALIUSPlacedTextContent{90.951}{111.928}{418.085}{ALIUSC000000}{\ALIUSFontCormorantRegular}{13.920}{making tasks satisfy similar conditions. When you are "forced" to decide at a}
  \ALIUSPlacedTextContent{90.951}{131.609}{418.215}{ALIUSC000000}{\ALIUSFontCormorantRegular}{13.920}{fixed cue, this can be thought of as driving a very strong urgency signal to}
  \ALIUSPlacedTextContent{90.951}{151.049}{418.299}{ALIUSC000000}{\ALIUSFontCormorantRegular}{13.920}{respond, similar to what I described previously. In this light, perhaps we}
  \ALIUSPlacedTextContent{90.951}{170.488}{418.331}{ALIUSC000000}{\ALIUSFontCormorantRegular}{13.920}{should not have been surprised to find the mechanism for changing the speed}
  \ALIUSPlacedTextContent{90.951}{189.928}{418.143}{ALIUSC000000}{\ALIUSFontCormorantRegular}{13.920}{vs. accuracy tradeoff that we discussed above involved boosting neural}
  \ALIUSPlacedTextContent{90.951}{209.609}{418.252}{ALIUSC000000}{\ALIUSFontCormorantRegular}{13.920}{responses rather than reducing the level of neural activity needed for}
  \ALIUSPlacedTextContent{90.951}{229.049}{414.974}{ALIUSC000000}{\ALIUSFontCormorantRegular}{13.920}{commitment, because no such reduction is seen for cued responses in forced-}
  \ALIUSPlacedTextContent{90.951}{248.489}{418.224}{ALIUSC000000}{\ALIUSFontCormorantRegular}{13.920}{choice tasks. Likewise, when a cue to respond comes relatively late, decision}
  \ALIUSPlacedTextContent{90.951}{267.928}{418.231}{ALIUSC000000}{\ALIUSFontCormorantRegular}{13.920}{commitment may have occurred already (Kiani et al., 2008). One of the}
  \ALIUSPlacedTextContent{90.951}{287.608}{418.195}{ALIUSC000000}{\ALIUSFontCormorantRegular}{13.920}{advantages of the cognitive nature of decision making is that it does not have}
  \ALIUSPlacedTextContent{90.951}{307.048}{418.257}{ALIUSC000000}{\ALIUSFontCormorantRegular}{13.920}{to be tied directly to the outside world, what one of my former advisors (Mike}
  \ALIUSPlacedTextContent{90.951}{326.488}{335.106}{ALIUSC000000}{\ALIUSFontCormorantRegular}{13.920}{Shadlen) would often describe as "freedom from immediacy".}
  \ALIUSPlacedTextContent{90.951}{358.053}{418.399}{ALIUSC1F8135}{\ALIUSFontLatoLight}{12.357}{Although it would stand to reason that the dynamics of the environment}
  \ALIUSPlacedTextContent{90.951}{375.333}{418.363}{ALIUSC1F8135}{\ALIUSFontLatoLight}{12.357}{dictate the optimal timescale (Ossmy et al., 2013; Glaze et al., 2015; Radillo}
  \ALIUSPlacedTextContent{90.951}{392.613}{418.463}{ALIUSC1F8135}{\ALIUSFontLatoLight}{12.357}{et al., 2017; Piet et al., 2018), you find that the brain represents and utilizes}
  \ALIUSPlacedTextContent{90.951}{409.893}{418.390}{ALIUSC1F8135}{\ALIUSFontLatoLight}{12.357}{multiple timescales of evidence evaluation during deliberation (Ganpuru et al.,}
  \ALIUSPlacedTextContent{90.951}{427.173}{418.382}{ALIUSC1F8135}{\ALIUSFontLatoLight}{12.357}{2019). What implications does this finding have for cellular mechanisms of}
  \ALIUSPlacedTextContent{90.951}{444.453}{418.413}{ALIUSC1F8135}{\ALIUSFontLatoLight}{12.357}{decision making? Does this imply a brain region involved in the task can have}
  \ALIUSPlacedTextContent{90.951}{461.733}{418.315}{ALIUSC1F8135}{\ALIUSFontLatoLight}{12.357}{heterogenous evidence weighting, or that multiple brain regions are}
  \ALIUSPlacedTextContent{90.951}{478.773}{232.665}{ALIUSC1F8135}{\ALIUSFontLatoLight}{12.357}{responsible for heterogeneous timescales?}
  \ALIUSPlacedTextContent{90.951}{507.928}{418.313}{ALIUSC000000}{\ALIUSFontCormorantRegular}{13.920}{We believe this speaks to the architecture of neural circuits that support}
  \ALIUSPlacedTextContent{90.951}{527.609}{418.287}{ALIUSC000000}{\ALIUSFontCormorantRegular}{13.920}{decision making. It establishes minimum capacities for information}
  \ALIUSPlacedTextContent{90.951}{547.049}{418.222}{ALIUSC000000}{\ALIUSFontCormorantRegular}{13.920}{processing that any neural mechanism of evidence evaluation must support.}
  \ALIUSPlacedTextContent{90.951}{566.489}{418.174}{ALIUSC000000}{\ALIUSFontCormorantRegular}{13.920}{Many existing models of decision processing have relied on precisely tuned}
  \ALIUSPlacedTextContent{90.951}{585.928}{418.269}{ALIUSC000000}{\ALIUSFontCormorantRegular}{13.920}{circuits that can effectively evaluate evidence over a set timescale (Wong \&}
  \ALIUSPlacedTextContent{90.951}{605.609}{418.347}{ALIUSC000000}{\ALIUSFontCormorantRegular}{13.920}{Wang, 2006). Changing the timescales can often be achieved by altering the}
  \ALIUSPlacedTextContent{90.951}{625.048}{418.367}{ALIUSC000000}{\ALIUSFontCormorantRegular}{13.920}{tuning of these circuits, but that is not enough to support multiple timescales}
  \ALIUSPlacedTextContent{90.951}{644.489}{418.364}{ALIUSC000000}{\ALIUSFontCormorantRegular}{13.920}{simultaneously. Intriguingly, alternative models already exist that can}
  \ALIUSPlacedTextContent{90.951}{663.928}{418.243}{ALIUSC000000}{\ALIUSFontCormorantRegular}{13.920}{support this. One such class of models uses a processing cascade with}
  \ALIUSPlacedTextContent{90.951}{683.369}{418.288}{ALIUSC000000}{\ALIUSFontCormorantRegular}{13.920}{progressively longer timescales of evaluation for later nodes in the network}
  \ALIUSPlacedTextContent{90.951}{703.048}{418.297}{ALIUSC000000}{\ALIUSFontCormorantRegular}{13.920}{(Goldman, 2009; Scott et al., 2017). Under this scheme, different timescales}
  \ALIUSPlacedTextContent{90.951}{722.489}{418.242}{ALIUSC000000}{\ALIUSFontCormorantRegular}{13.920}{are represented by different network nodes without requiring any alterations}
  \ALIUSPlacedTextContent{90.951}{741.929}{418.237}{ALIUSC000000}{\ALIUSFontCormorantRegular}{13.920}{in tuning, so it could naturally support simultaneous deliberation across}
  \ALIUSPlacedTextContent{90.951}{793.143}{122.397}{ALIUSC767171}{\ALIUSFontCormorantLight}{12.000}{ALIUS Bulletin n°4 (2020)}
  \ALIUSPlacedTextContent{387.435}{793.143}{120.728}{ALIUSC767171}{\ALIUSFontCormorantLight}{12.000}{aliusresearch.org/bulletin}
\end{tikzpicture}
\null
\newpage
% Page 9
\thispagestyle{empty}
\noindent\begin{tikzpicture}[remember picture,overlay,x=1bp,y=1bp,shift={(current page.north west)}]
  \ALIUSPlacedTextContent{90.951}{36.329}{321.512}{ALIUSC767171}{\ALIUSFontCormorantRegular}{13.920}{Timothy Hanks – The Neuroscience of Decision Making}
  \ALIUSPlacedTextContent{493.431}{36.329}{15.773}{ALIUSC000000}{\ALIUSFontCormorantRegular}{13.920}{76}
  \ALIUSPlacedTextContent{90.951}{73.048}{418.024}{ALIUSC000000}{\ALIUSFontCormorantRegular}{13.920}{multiple timescales. We have not yet shown this to be the case, but we think}
  \ALIUSPlacedTextContent{90.951}{92.488}{418.240}{ALIUSC000000}{\ALIUSFontCormorantRegular}{13.920}{it is a good hypothesis. In theory, these nodes could be distributed across}
  \ALIUSPlacedTextContent{90.951}{111.928}{418.344}{ALIUSC000000}{\ALIUSFontCormorantRegular}{13.920}{different brain regions or be distinct subcomponents within brain regions}
  \ALIUSPlacedTextContent{90.951}{131.609}{385.841}{ALIUSC000000}{\ALIUSFontCormorantRegular}{13.920}{related to decision making, and this is also a question we are pursuing.}
  \ALIUSPlacedTextContent{90.951}{163.173}{418.469}{ALIUSC1F8135}{\ALIUSFontLatoLight}{12.357}{Disease states can influence evidence accumulation, resulting in too much or}
  \ALIUSPlacedTextContent{90.951}{180.453}{418.375}{ALIUSC1F8135}{\ALIUSFontLatoLight}{12.357}{too little exploration in various domains. How is decision making affected in}
  \ALIUSPlacedTextContent{90.951}{197.733}{418.357}{ALIUSC1F8135}{\ALIUSFontLatoLight}{12.357}{schizophrenia and depression? Does this point to underlying mechanistic}
  \ALIUSPlacedTextContent{90.951}{215.013}{418.517}{ALIUSC1F8135}{\ALIUSFontLatoLight}{12.357}{etiological similarities among disease states, or computational attractor states}
  \ALIUSPlacedTextContent{90.951}{232.053}{418.373}{ALIUSC1F8135}{\ALIUSFontLatoLight}{12.357}{that can be induced by a variety of causes? What are promising cellular}
  \ALIUSPlacedTextContent{90.951}{249.333}{414.986}{ALIUSC1F8135}{\ALIUSFontLatoLight}{12.357}{mechanisms or affected circuits to investigate in relationship to decision-}
  \ALIUSPlacedTextContent{90.951}{266.613}{89.099}{ALIUSC1F8135}{\ALIUSFontLatoLight}{12.357}{making deficits?}
  \ALIUSPlacedTextContent{90.951}{295.768}{418.355}{ALIUSC000000}{\ALIUSFontCormorantRegular}{13.920}{Impaired decision making is observed in almost all mental disorders,}
  \ALIUSPlacedTextContent{90.951}{315.208}{418.246}{ALIUSC000000}{\ALIUSFontCormorantRegular}{13.920}{including schizophrenia and depression, as you mention. Furthermore, these}
  \ALIUSPlacedTextContent{90.951}{334.888}{418.356}{ALIUSC000000}{\ALIUSFontCormorantRegular}{13.920}{impairments can often severely impact quality of life, and some of the largest}
  \ALIUSPlacedTextContent{90.951}{354.328}{418.257}{ALIUSC000000}{\ALIUSFontCormorantRegular}{13.920}{negative impacts from many mental disorders come from poor decision}
  \ALIUSPlacedTextContent{90.951}{373.768}{418.135}{ALIUSC000000}{\ALIUSFontCormorantRegular}{13.920}{making. I saw this firsthand with my dad, who suffered from an atypical}
  \ALIUSPlacedTextContent{90.951}{393.208}{418.247}{ALIUSC000000}{\ALIUSFontCormorantRegular}{13.920}{variant of Alzheimer's disease that presented as frontal-temporal dementia.}
  \ALIUSPlacedTextContent{90.951}{412.888}{418.229}{ALIUSC000000}{\ALIUSFontCormorantRegular}{13.920}{The biggest impacts on his life early on in the disease progression were with}
  \ALIUSPlacedTextContent{90.951}{432.328}{94.955}{ALIUSC000000}{\ALIUSFontCormorantRegular}{13.920}{decision making.}
  \ALIUSPlacedTextContent{90.951}{463.768}{418.253}{ALIUSC000000}{\ALIUSFontCormorantRegular}{13.920}{While decision making can be impaired in a variety of ways with different}
  \ALIUSPlacedTextContent{90.951}{483.208}{418.253}{ALIUSC000000}{\ALIUSFontCormorantRegular}{13.920}{disorders, there do seem to be common modes of impairment that can occur}
  \ALIUSPlacedTextContent{90.951}{502.648}{418.184}{ALIUSC000000}{\ALIUSFontCormorantRegular}{13.920}{across disorders. What we are trying to do now is exactly in line with what}
  \ALIUSPlacedTextContent{90.951}{522.328}{418.157}{ALIUSC000000}{\ALIUSFontCormorantRegular}{13.920}{you propose of applying our computational models to describe these modes}
  \ALIUSPlacedTextContent{90.951}{541.768}{418.332}{ALIUSC000000}{\ALIUSFontCormorantRegular}{13.920}{of impairment and link them to underlying neural mechanisms. We have}
  \ALIUSPlacedTextContent{90.951}{561.208}{418.256}{ALIUSC000000}{\ALIUSFontCormorantRegular}{13.920}{begun to collaborate with UC Davis Conte Center led by Drs. Cameron}
  \ALIUSPlacedTextContent{90.951}{580.648}{418.179}{ALIUSC000000}{\ALIUSFontCormorantRegular}{13.920}{Carter and Kimberly McAllister to apply this to schizophrenia patients and}
  \ALIUSPlacedTextContent{90.951}{600.328}{418.262}{ALIUSC000000}{\ALIUSFontCormorantRegular}{13.920}{animal models aimed at figuring out neural mechanisms involved. We believe}
  \ALIUSPlacedTextContent{90.951}{619.768}{418.259}{ALIUSC000000}{\ALIUSFontCormorantRegular}{13.920}{that corticostriatal circuits may play a central role in the impairments to}
  \ALIUSPlacedTextContent{90.951}{639.208}{418.274}{ALIUSC000000}{\ALIUSFontCormorantRegular}{13.920}{decision making that come about with schizophrenia. With Dr. Randy}
  \ALIUSPlacedTextContent{90.951}{658.648}{418.176}{ALIUSC000000}{\ALIUSFontCormorantRegular}{13.920}{O'Reilly, we have developed a computational model for impaired decision}
  \ALIUSPlacedTextContent{90.951}{678.328}{418.204}{ALIUSC000000}{\ALIUSFontCormorantRegular}{13.920}{making in schizophrenia that makes specific predictions for neural}
  \ALIUSPlacedTextContent{90.951}{697.768}{418.284}{ALIUSC000000}{\ALIUSFontCormorantRegular}{13.920}{mechanisms impacted in corticostriatal circuits, and we are beginning to test}
  \ALIUSPlacedTextContent{90.951}{717.208}{100.093}{ALIUSC000000}{\ALIUSFontCormorantRegular}{13.920}{those predictions.}
  \ALIUSPlacedTextContent{90.951}{793.143}{122.397}{ALIUSC767171}{\ALIUSFontCormorantLight}{12.000}{ALIUS Bulletin n°4 (2020)}
  \ALIUSPlacedTextContent{387.435}{793.143}{120.728}{ALIUSC767171}{\ALIUSFontCormorantLight}{12.000}{aliusresearch.org/bulletin}
\end{tikzpicture}
\null
\newpage
% Page 10
\thispagestyle{empty}
\noindent\begin{tikzpicture}[remember picture,overlay,x=1bp,y=1bp,shift={(current page.north west)}]
  \ALIUSPlacedTextContent{90.951}{36.329}{321.512}{ALIUSC767171}{\ALIUSFontCormorantRegular}{13.920}{Timothy Hanks – The Neuroscience of Decision Making}
  \ALIUSPlacedTextContent{493.911}{36.329}{15.269}{ALIUSC000000}{\ALIUSFontCormorantRegular}{13.920}{77}
  \ALIUSPlacedTextContent{98.191}{76.599}{14.160}{ALIUSC595959}{\ALIUSFontCormorantLight}{48.000}{\ALIUSPullQuoteOpen}
  \ALIUSPlacedTextContent{133.951}{89.947}{322.721}{ALIUSC000000}{\ALIUSFontLatoLight}{15.120}{Attentional allocation can be viewed as a decision}
  \ALIUSPlacedTextContent{133.951}{107.947}{322.875}{ALIUSC000000}{\ALIUSFontLatoLight}{15.120}{process, and a reduced bound that would explain}
  \ALIUSPlacedTextContent{133.951}{125.947}{322.742}{ALIUSC000000}{\ALIUSFontLatoLight}{15.120}{impulsivity would also explain lower thresholds}
  \ALIUSPlacedTextContent{469.951}{140.679}{14.160}{ALIUSC595959}{\ALIUSFontCormorantLight}{48.000}{\ALIUSPullQuoteClose}
  \ALIUSPlacedTextContent{133.951}{143.947}{193.377}{ALIUSC000000}{\ALIUSFontLatoLight}{15.120}{for changing attentional focus}
  \ALIUSPlacedTextContent{90.951}{193.049}{418.296}{ALIUSC000000}{\ALIUSFontCormorantRegular}{13.920}{We have also been collaborating with Dr. Johannes Hell to study neural}
  \ALIUSPlacedTextContent{90.951}{212.488}{418.154}{ALIUSC000000}{\ALIUSFontCormorantRegular}{13.920}{mechanisms of attentional control that might be affected in ADHD. One of}
  \ALIUSPlacedTextContent{90.951}{231.928}{418.300}{ALIUSC000000}{\ALIUSFontCormorantRegular}{13.920}{the most interesting aspects of ADHD is that it results in impairments to not}
  \ALIUSPlacedTextContent{90.951}{251.609}{418.280}{ALIUSC000000}{\ALIUSFontCormorantRegular}{13.920}{only attention, but it also leads to more impulsivity. What is the connection}
  \ALIUSPlacedTextContent{90.951}{271.048}{418.231}{ALIUSC000000}{\ALIUSFontCormorantRegular}{13.920}{between the two? Through our work on decision commitment, we have a}
  \ALIUSPlacedTextContent{90.951}{290.488}{418.232}{ALIUSC000000}{\ALIUSFontCormorantRegular}{13.920}{good handle on mechanisms that may be involved for impulsivity, and again,}
  \ALIUSPlacedTextContent{90.951}{309.928}{418.253}{ALIUSC000000}{\ALIUSFontCormorantRegular}{13.920}{corticostriatal circuits are implicated. I believe that similar circuit}
  \ALIUSPlacedTextContent{90.951}{329.608}{418.182}{ALIUSC000000}{\ALIUSFontCormorantRegular}{13.920}{mechanisms may also underlie attentional deficits. In particular, individuals}
  \ALIUSPlacedTextContent{90.951}{349.048}{418.235}{ALIUSC000000}{\ALIUSFontCormorantRegular}{13.920}{with ADHD are often able to reap the benefits of attention when they are}
  \ALIUSPlacedTextContent{90.951}{368.488}{418.271}{ALIUSC000000}{\ALIUSFontCormorantRegular}{13.920}{attending to something, but the problem is that they are more easily}
  \ALIUSPlacedTextContent{90.951}{387.928}{418.267}{ALIUSC000000}{\ALIUSFontCormorantRegular}{13.920}{distracted. It isn't so much a problem with attentional enhancement of}
  \ALIUSPlacedTextContent{90.951}{407.608}{418.172}{ALIUSC000000}{\ALIUSFontCormorantRegular}{13.920}{perception, but rather one of allocating those attentional resources.}
  \ALIUSPlacedTextContent{90.951}{427.048}{418.254}{ALIUSC000000}{\ALIUSFontCormorantRegular}{13.920}{Attentional allocation can be viewed as a decision process, and a reduced}
  \ALIUSPlacedTextContent{90.951}{446.488}{418.363}{ALIUSC000000}{\ALIUSFontCormorantRegular}{13.920}{bound that would explain impulsivity would also explain lower thresholds}
  \ALIUSPlacedTextContent{90.951}{465.928}{418.332}{ALIUSC000000}{\ALIUSFontCormorantRegular}{13.920}{for changing attentional focus – the distractibility that affects those with}
  \ALIUSPlacedTextContent{90.951}{485.608}{418.171}{ALIUSC000000}{\ALIUSFontCormorantRegular}{13.920}{ADHD. We think that corticostriatal circuit mechanisms may explain both.}
  \ALIUSPlacedTextContent{90.951}{505.048}{418.219}{ALIUSC000000}{\ALIUSFontCormorantRegular}{13.920}{With Dr. Hell, we are trying to determine specific cellular mechanisms that}
  \ALIUSPlacedTextContent{90.951}{524.489}{94.465}{ALIUSC000000}{\ALIUSFontCormorantRegular}{13.920}{may be involved.}
  \ALIUSPlacedTextContent{90.951}{556.053}{418.391}{ALIUSC1F8135}{\ALIUSFontLatoLight}{12.357}{What is the minimal number of neurons required to 'make a decision'? More}
  \ALIUSPlacedTextContent{90.951}{573.333}{418.522}{ALIUSC1F8135}{\ALIUSFontLatoLight}{12.357}{generally, do neurons as a cell type participate in a special type of decision}
  \ALIUSPlacedTextContent{90.951}{590.613}{418.424}{ALIUSC1F8135}{\ALIUSFontLatoLight}{12.357}{making relative to non-neuronal cells or perhaps even non-cells? In other}
  \ALIUSPlacedTextContent{90.951}{607.893}{418.383}{ALIUSC1F8135}{\ALIUSFontLatoLight}{12.357}{words, how should we study the "decision making" processes of bacteria,}
  \ALIUSPlacedTextContent{90.951}{625.173}{241.921}{ALIUSC1F8135}{\ALIUSFontLatoLight}{12.357}{computers, neural and non-neural cell types?}
  \ALIUSPlacedTextContent{90.951}{654.328}{418.441}{ALIUSC000000}{\ALIUSFontCormorantRegular}{13.920}{I take an inclusive view on what it means to "make a decision", but with the}
  \ALIUSPlacedTextContent{90.951}{673.768}{418.210}{ALIUSC000000}{\ALIUSFontCormorantRegular}{13.920}{caveat that there are important differences among the class of processes}
  \ALIUSPlacedTextContent{90.951}{693.208}{418.301}{ALIUSC000000}{\ALIUSFontCormorantRegular}{13.920}{considered as such. Under this view, a single neuron is certainly capable of}
  \ALIUSPlacedTextContent{90.951}{712.648}{418.374}{ALIUSC000000}{\ALIUSFontCormorantRegular}{13.920}{making a decision, but without the same range of complexity as a large}
  \ALIUSPlacedTextContent{90.951}{732.328}{418.213}{ALIUSC000000}{\ALIUSFontCormorantRegular}{13.920}{network of neurons. Likewise, while neurons and neural networks participate}
  \ALIUSPlacedTextContent{90.951}{751.768}{418.207}{ALIUSC000000}{\ALIUSFontCormorantRegular}{13.920}{in a variety of types of decision making, they do not do so in a necessarily}
  \ALIUSPlacedTextContent{90.951}{793.143}{122.397}{ALIUSC767171}{\ALIUSFontCormorantLight}{12.000}{ALIUS Bulletin n°4 (2020)}
  \ALIUSPlacedTextContent{387.435}{793.143}{120.728}{ALIUSC767171}{\ALIUSFontCormorantLight}{12.000}{aliusresearch.org/bulletin}
\end{tikzpicture}
\null
\newpage
% Page 11
\thispagestyle{empty}
\noindent\begin{tikzpicture}[remember picture,overlay,x=1bp,y=1bp,shift={(current page.north west)}]
  \fill[ALIUSC0000FF] (126.951bp,-749.031bp) rectangle ++(217.920bp,-0.720bp);
  \ALIUSPlacedTextContent{90.951}{36.329}{321.512}{ALIUSC767171}{\ALIUSFontCormorantRegular}{13.920}{Timothy Hanks – The Neuroscience of Decision Making}
  \ALIUSPlacedTextContent{493.191}{36.329}{16.109}{ALIUSC000000}{\ALIUSFontCormorantRegular}{13.920}{78}
  \ALIUSPlacedTextContent{90.951}{73.048}{418.261}{ALIUSC000000}{\ALIUSFontCormorantRegular}{13.920}{unique way. Non-neural cell types, computers, etc. may also participate in}
  \ALIUSPlacedTextContent{90.951}{92.488}{414.974}{ALIUSC000000}{\ALIUSFontCormorantRegular}{13.920}{similar types of decision making, but this needs to be considered on a case-}
  \ALIUSPlacedTextContent{90.951}{111.928}{418.222}{ALIUSC000000}{\ALIUSFontCormorantRegular}{13.920}{by-case basis. For instance, we can describe accumulating information to a}
  \ALIUSPlacedTextContent{90.951}{131.609}{418.410}{ALIUSC000000}{\ALIUSFontCormorantRegular}{13.920}{threshold over a flexible timescale from milliseconds to minutes as a specific}
  \ALIUSPlacedTextContent{90.951}{151.049}{418.409}{ALIUSC000000}{\ALIUSFontCormorantRegular}{13.920}{type of decision making. Human brains can do this, but it is not known what}
  \ALIUSPlacedTextContent{90.951}{170.488}{418.165}{ALIUSC000000}{\ALIUSFontCormorantRegular}{13.920}{minimal architecture of brain tissue is needed to accomplish it. Computers}
  \ALIUSPlacedTextContent{90.951}{189.928}{418.279}{ALIUSC000000}{\ALIUSFontCormorantRegular}{13.920}{can easily be programmed to do it, but it becomes more challenging as we}
  \ALIUSPlacedTextContent{90.951}{209.609}{418.234}{ALIUSC000000}{\ALIUSFontCormorantRegular}{13.920}{expand the domain that we are considering to more complex types of}
  \ALIUSPlacedTextContent{90.951}{229.049}{418.233}{ALIUSC000000}{\ALIUSFontCormorantRegular}{13.920}{decision making. It would also be interesting to know the capacity of other}
  \ALIUSPlacedTextContent{90.951}{248.489}{418.350}{ALIUSC000000}{\ALIUSFontCormorantRegular}{13.920}{systems like bacteria. Intriguingly, for this example, there are suggestions that}
  \ALIUSPlacedTextContent{90.951}{267.928}{414.974}{ALIUSC000000}{\ALIUSFontCormorantRegular}{13.920}{non-neuronal cells could play a role. In particular, the glial cells (non-}
  \ALIUSPlacedTextContent{90.951}{287.608}{418.399}{ALIUSC000000}{\ALIUSFontCormorantRegular}{13.920}{neuronal cells of the nervous system that actually outnumber neurons)}
  \ALIUSPlacedTextContent{90.951}{307.048}{418.302}{ALIUSC000000}{\ALIUSFontCormorantRegular}{13.920}{known as astrocytes have been shown to accumulate information over 10s of}
  \ALIUSPlacedTextContent{90.951}{326.488}{418.355}{ALIUSC000000}{\ALIUSFontCormorantRegular}{13.920}{seconds for decision making in zebrafish, with some indications of flexibility}
  \ALIUSPlacedTextContent{90.951}{345.928}{418.238}{ALIUSC000000}{\ALIUSFontCormorantRegular}{13.920}{in that timescale. In that circuit, the astrocytes seem to work in concert with}
  \ALIUSPlacedTextContent{90.951}{365.608}{418.157}{ALIUSC000000}{\ALIUSFontCormorantRegular}{13.920}{neurons by accumulating evidence from neural input over a longer timescale}
  \ALIUSPlacedTextContent{90.951}{385.048}{198.749}{ALIUSC000000}{\ALIUSFontCormorantRegular}{13.920}{and then influencing other neurons.}
  \ALIUSPlacedTextContent{90.951}{416.488}{418.309}{ALIUSC000000}{\ALIUSFontCormorantRegular}{13.920}{I believe one of the keys for studying decision making in any setting is striving}
  \ALIUSPlacedTextContent{90.951}{435.928}{418.178}{ALIUSC000000}{\ALIUSFontCormorantRegular}{13.920}{for a clear description of what type of decision process one is studying. This}
  \ALIUSPlacedTextContent{90.951}{455.608}{418.276}{ALIUSC000000}{\ALIUSFontCormorantRegular}{13.920}{is one of the more important benefits of using a computational approach}
  \ALIUSPlacedTextContent{90.951}{475.048}{418.203}{ALIUSC000000}{\ALIUSFontCormorantRegular}{13.920}{because it gives a framework for rigorous description. This has a parallel to}
  \ALIUSPlacedTextContent{90.951}{494.488}{418.303}{ALIUSC000000}{\ALIUSFontCormorantRegular}{13.920}{the discussions above about free will. Rather than asking, "Is [X] making a}
  \ALIUSPlacedTextContent{90.951}{513.928}{418.282}{ALIUSC000000}{\ALIUSFontCormorantRegular}{13.920}{decision?", I think it is better to ask, "In what way is [X] making a decision?"}
  \ALIUSPlacedTextContent{90.951}{533.369}{418.266}{ALIUSC000000}{\ALIUSFontCormorantRegular}{13.920}{Reflexive movements may still be considered decisions in some sense, but}
  \ALIUSPlacedTextContent{90.951}{553.049}{418.261}{ALIUSC000000}{\ALIUSFontCormorantRegular}{13.920}{they are very different types of decisions than those involving reflective}
  \ALIUSPlacedTextContent{90.951}{572.489}{418.226}{ALIUSC000000}{\ALIUSFontCormorantRegular}{13.920}{deliberation. They invoke different forms of computation and different}
  \ALIUSPlacedTextContent{90.951}{591.928}{217.984}{ALIUSC000000}{\ALIUSFontCormorantRegular}{13.920}{complexity of neural circuit processing.}
  \ALIUSPlacedTextContent{264.617}{643.132}{71.454}{ALIUSC000000}{\ALIUSFontLatoLight}{13.920}{References}
  \ALIUSPlacedTextContent{90.951}{688.815}{418.030}{ALIUSC000000}{\ALIUSFontCormorantRegular}{12.960}{Beck J. M., Ma W. J., Kiani R., Hanks T., Churchland A. K., Roitman J., Shadlen,}
  \ALIUSPlacedTextContent{126.951}{704.656}{382.046}{ALIUSC000000}{\ALIUSFontCormorantRegular}{12.960}{M. N., \& Pouget, A. (2008) Probabilistic population codes for Bayesian}
  \ALIUSPlacedTextContent{126.951}{720.256}{44.774}{ALIUSC000000}{\ALIUSFontCormorantRegular}{12.960}{decision}
  \ALIUSPlacedTextContent{218.077}{720.256}{43.406}{ALIUSC000000}{\ALIUSFontCormorantRegular}{12.960}{making.}
  \ALIUSPlacedTextContent{307.825}{720.256}{34.577}{ALIUSC000000}{\ALIUSFontCormorantItalic}{12.960}{Neuron}
  \ALIUSPlacedTextContent{342.431}{720.256}{5.919}{ALIUSC000000}{\ALIUSFontCormorantRegular}{12.960}{,}
  \ALIUSPlacedTextContent{394.701}{720.256}{13.707}{ALIUSC000000}{\ALIUSFontCormorantItalic}{12.960}{60.}
  \ALIUSPlacedTextContent{457.800}{720.256}{51.183}{ALIUSC000000}{\ALIUSFontCormorantRegular}{12.960}{1142–1152.}
  \ALIUSPlacedTextContent{126.951}{736.096}{218.030}{ALIUSC0000FF}{\ALIUSFontCormorantRegular}{12.960}{https://doi.org/10.1016/j.neuron.2008.09.021}
  \ALIUSPlacedTextContent{90.951}{793.143}{122.397}{ALIUSC767171}{\ALIUSFontCormorantLight}{12.000}{ALIUS Bulletin n°4 (2020)}
  \ALIUSPlacedTextContent{387.435}{793.143}{120.728}{ALIUSC767171}{\ALIUSFontCormorantLight}{12.000}{aliusresearch.org/bulletin}
\end{tikzpicture}
\null
\newpage
% Page 12
\thispagestyle{empty}
\noindent\begin{tikzpicture}[remember picture,overlay,x=1bp,y=1bp,shift={(current page.north west)}]
  \fill[ALIUSC0000FF] (126.951bp,-117.351bp) rectangle ++(204.480bp,-0.720bp);
  \fill[ALIUSC0000FF] (197.751bp,-180.231bp) rectangle ++(198.480bp,-0.720bp);
  \fill[ALIUSC0000FF] (300.471bp,-227.511bp) rectangle ++(174.480bp,-0.720bp);
  \fill[ALIUSC0000FF] (126.951bp,-285.831bp) rectangle ++(272.640bp,-0.720bp);
  \fill[ALIUSC0000FF] (171.831bp,-332.871bp) rectangle ++(215.280bp,-0.720bp);
  \fill[ALIUSC0000FF] (316.791bp,-380.151bp) rectangle ++(174.000bp,-0.720bp);
  \fill[ALIUSC0000FF] (126.951bp,-443.271bp) rectangle ++(214.800bp,-0.720bp);
  \fill[ALIUSC0000FF] (229.671bp,-490.311bp) rectangle ++(216.960bp,-0.720bp);
  \fill[ALIUSC0000FF] (126.951bp,-569.031bp) rectangle ++(179.520bp,-0.720bp);
  \fill[ALIUSC0000FF] (126.951bp,-647.751bp) rectangle ++(250.320bp,-0.720bp);
  \fill[ALIUSC0000FF] (126.951bp,-726.471bp) rectangle ++(198.480bp,-0.720bp);
  \ALIUSPlacedTextContent{90.951}{36.329}{321.512}{ALIUSC767171}{\ALIUSFontCormorantRegular}{13.920}{Timothy Hanks – The Neuroscience of Decision Making}
  \ALIUSPlacedTextContent{493.431}{36.329}{15.773}{ALIUSC000000}{\ALIUSFontCormorantRegular}{13.920}{79}
  \ALIUSPlacedTextContent{90.951}{72.976}{418.009}{ALIUSC000000}{\ALIUSFontCormorantRegular}{12.960}{Brody C. D., Hanks T. D. (2016). Neural underpinnings of the evidence}
  \ALIUSPlacedTextContent{126.951}{88.576}{68.605}{ALIUSC000000}{\ALIUSFontCormorantRegular}{12.960}{accumulator.}
  \ALIUSPlacedTextContent{195.576}{88.576}{141.875}{ALIUSC000000}{\ALIUSFontCormorantItalic}{12.960}{Current Opinion Neurobiology}
  \ALIUSPlacedTextContent{337.468}{88.576}{5.593}{ALIUSC000000}{\ALIUSFontCormorantRegular}{12.960}{.}
  \ALIUSPlacedTextContent{343.070}{88.576}{9.461}{ALIUSC000000}{\ALIUSFontCormorantItalic}{12.960}{37}
  \ALIUSPlacedTextContent{352.547}{88.576}{49.324}{ALIUSC000000}{\ALIUSFontCormorantRegular}{12.960}{. 149–157.}
  \ALIUSPlacedTextContent{126.951}{104.416}{204.616}{ALIUSC0000FF}{\ALIUSFontCormorantRegular}{12.960}{https://doi.org/10.1016/j.conb.2016.01.003}
  \ALIUSPlacedTextContent{90.951}{135.856}{418.023}{ALIUSC000000}{\ALIUSFontCormorantRegular}{12.960}{Ganupuru P., Goldring A. B., Harun R., Hanks T. D. (2019). Flexibility of}
  \ALIUSPlacedTextContent{126.951}{151.696}{286.611}{ALIUSC000000}{\ALIUSFontCormorantRegular}{12.960}{Timescales of Evidence Evaluation for Decision Making.}
  \ALIUSPlacedTextContent{413.634}{151.696}{89.739}{ALIUSC000000}{\ALIUSFontCormorantItalic}{12.960}{Current Biology. 29}
  \ALIUSPlacedTextContent{503.390}{151.696}{5.594}{ALIUSC000000}{\ALIUSFontCormorantRegular}{12.960}{:}
  \ALIUSPlacedTextContent{126.951}{167.296}{70.831}{ALIUSC000000}{\ALIUSFontCormorantRegular}{12.960}{2091–2097.e4.}
  \ALIUSPlacedTextContent{197.812}{167.296}{198.434}{ALIUSC0000FF}{\ALIUSFontCormorantRegular}{12.960}{https://doi.org/10.1016/j.cub.2019.05.037}
  \ALIUSPlacedTextContent{90.951}{198.736}{400.575}{ALIUSC000000}{\ALIUSFontCormorantRegular}{12.960}{Glaze C. M., Kable J. W., Gold J. I. (2015). Normative evidence accumulation in}
  \ALIUSPlacedTextContent{112.951}{214.576}{148.735}{ALIUSC000000}{\ALIUSFontCormorantRegular}{12.960}{unpredictable environments.}
  \ALIUSPlacedTextContent{261.706}{214.576}{21.969}{ALIUSC000000}{\ALIUSFontCormorantItalic}{12.960}{eLife}
  \ALIUSPlacedTextContent{283.689}{214.576}{5.593}{ALIUSC000000}{\ALIUSFontCormorantRegular}{12.960}{.}
  \ALIUSPlacedTextContent{289.292}{214.576}{5.482}{ALIUSC000000}{\ALIUSFontCormorantItalic}{12.960}{4}
  \ALIUSPlacedTextContent{294.791}{214.576}{5.594}{ALIUSC000000}{\ALIUSFontCormorantRegular}{12.960}{.}
  \ALIUSPlacedTextContent{300.393}{214.576}{174.562}{ALIUSC0000FF}{\ALIUSFontCormorantRegular}{12.960}{https://doi.org/10.7554/eLife.08825}
  \ALIUSPlacedTextContent{90.951}{241.216}{344.430}{ALIUSC000000}{\ALIUSFontCormorantRegular}{12.960}{Gold J. I., Shadlen M. N. (2007). The neural basis of decision making.}
  \ALIUSPlacedTextContent{434.943}{241.216}{74.036}{ALIUSC000000}{\ALIUSFontCormorantItalic}{12.960}{Annual Review}
  \ALIUSPlacedTextContent{126.951}{257.056}{75.354}{ALIUSC000000}{\ALIUSFontCormorantItalic}{12.960}{Neuroscience. 30}
  \ALIUSPlacedTextContent{202.323}{257.056}{50.338}{ALIUSC000000}{\ALIUSFontCormorantRegular}{12.960}{. 535–574.}
  \ALIUSPlacedTextContent{126.951}{272.896}{272.553}{ALIUSC0000FF}{\ALIUSFontCormorantRegular}{12.960}{https://doi.org/10.1146/annurev.neuro.29.051605.113038}
  \ALIUSPlacedTextContent{90.951}{304.336}{362.273}{ALIUSC000000}{\ALIUSFontCormorantRegular}{12.960}{Goldman M. S. (2009). Memory without feedback in a neural network.}
  \ALIUSPlacedTextContent{453.552}{304.336}{49.818}{ALIUSC000000}{\ALIUSFontCormorantItalic}{12.960}{Neuron. 61}
  \ALIUSPlacedTextContent{503.390}{304.336}{5.593}{ALIUSC000000}{\ALIUSFontCormorantRegular}{12.960}{.}
  \ALIUSPlacedTextContent{126.951}{319.936}{44.840}{ALIUSC000000}{\ALIUSFontCormorantRegular}{12.960}{621–634.}
  \ALIUSPlacedTextContent{171.799}{319.936}{215.321}{ALIUSC0000FF}{\ALIUSFontCormorantRegular}{12.960}{https://doi.org/10.1016/j.neuron.2008.12.012}
  \ALIUSPlacedTextContent{90.951}{351.616}{418.034}{ALIUSC000000}{\ALIUSFontCormorantRegular}{12.960}{Hanks T., Kiani R., Shadlen M. N. (2014). A neural mechanism of speed-accuracy}
  \ALIUSPlacedTextContent{126.951}{367.216}{152.128}{ALIUSC000000}{\ALIUSFontCormorantRegular}{12.960}{tradeoff in macaque area LIP.}
  \ALIUSPlacedTextContent{279.125}{367.216}{32.018}{ALIUSC000000}{\ALIUSFontCormorantItalic}{12.960}{eLife. 3}
  \ALIUSPlacedTextContent{311.156}{367.216}{5.594}{ALIUSC000000}{\ALIUSFontCormorantRegular}{12.960}{.}
  \ALIUSPlacedTextContent{316.759}{367.216}{173.979}{ALIUSC0000FF}{\ALIUSFontCormorantRegular}{12.960}{https://doi.org/10.7554/eLife.02260}
  \ALIUSPlacedTextContent{90.951}{398.656}{418.075}{ALIUSC000000}{\ALIUSFontCormorantRegular}{12.960}{Hanks T. D., Summerfield C. (2017). Perceptual Decision Making in Rodents,}
  \ALIUSPlacedTextContent{126.951}{414.496}{121.707}{ALIUSC000000}{\ALIUSFontCormorantRegular}{12.960}{Monkeys, and Humans.}
  \ALIUSPlacedTextContent{248.654}{414.496}{50.101}{ALIUSC000000}{\ALIUSFontCormorantItalic}{12.960}{Neuron. 93}
  \ALIUSPlacedTextContent{298.768}{414.496}{47.215}{ALIUSC000000}{\ALIUSFontCormorantRegular}{12.960}{. 15–31.}
  \ALIUSPlacedTextContent{126.951}{430.336}{214.689}{ALIUSC0000FF}{\ALIUSFontCormorantRegular}{12.960}{https://doi.org/10.1016/j.neuron.2016.12.003}
  \ALIUSPlacedTextContent{90.951}{461.776}{418.043}{ALIUSC000000}{\ALIUSFontCormorantRegular}{12.960}{Heitz R. P., Schall J. D. (2012). Neural mechanisms of speed-accuracy tradeoff.}
  \ALIUSPlacedTextContent{126.951}{477.376}{50.674}{ALIUSC000000}{\ALIUSFontCormorantItalic}{12.960}{Neuron. 76}
  \ALIUSPlacedTextContent{177.650}{477.376}{51.923}{ALIUSC000000}{\ALIUSFontCormorantRegular}{12.960}{. 616–628.}
  \ALIUSPlacedTextContent{229.583}{477.376}{217.032}{ALIUSC0000FF}{\ALIUSFontCormorantRegular}{12.960}{https://doi.org/10.1016/j.neuron.2012.08.030}
  \ALIUSPlacedTextContent{90.951}{509.056}{418.074}{ALIUSC000000}{\ALIUSFontCormorantRegular}{12.960}{Johnson B., Verma R., Sun M., Hanks T. D. (2017). Characterization of decision}
  \ALIUSPlacedTextContent{126.951}{524.656}{382.054}{ALIUSC000000}{\ALIUSFontCormorantRegular}{12.960}{commitment rule alterations during an auditory change detection task.}
  \ALIUSPlacedTextContent{126.951}{540.496}{144.227}{ALIUSC000000}{\ALIUSFontCormorantItalic}{12.960}{Journal of Neurophysiology. 118}
  \ALIUSPlacedTextContent{271.195}{540.496}{61.349}{ALIUSC000000}{\ALIUSFontCormorantRegular}{12.960}{. 2526–2536.}
  \ALIUSPlacedTextContent{126.951}{556.096}{179.616}{ALIUSC0000FF}{\ALIUSFontCormorantRegular}{12.960}{https://doi.org/10.1152/jn.00071.2017}
  \ALIUSPlacedTextContent{90.951}{587.536}{418.071}{ALIUSC000000}{\ALIUSFontCormorantRegular}{12.960}{Kiani R., Hanks T. D., Shadlen M. N. (2008). Bounded integration in parietal cortex}
  \ALIUSPlacedTextContent{126.951}{603.376}{382.029}{ALIUSC000000}{\ALIUSFontCormorantRegular}{12.960}{underlies decisions even when viewing duration is dictated by the}
  \ALIUSPlacedTextContent{126.951}{619.216}{70.808}{ALIUSC000000}{\ALIUSFontCormorantRegular}{12.960}{environment.}
  \ALIUSPlacedTextContent{197.786}{619.216}{124.364}{ALIUSC000000}{\ALIUSFontCormorantItalic}{12.960}{Journal of Neuroscience. 28}
  \ALIUSPlacedTextContent{322.168}{619.216}{59.010}{ALIUSC000000}{\ALIUSFontCormorantRegular}{12.960}{. 3017–3029}
  \ALIUSPlacedTextContent{126.951}{634.815}{250.217}{ALIUSC0000FF}{\ALIUSFontCormorantRegular}{12.960}{https://doi.org/10.1523/JNEUROSCI.4761-07.2008}
  \ALIUSPlacedTextContent{90.951}{666.256}{418.023}{ALIUSC000000}{\ALIUSFontCormorantRegular}{12.960}{Ossmy O., Moran R., Pfeffer T., Tsetsos K., Usher M., Donner T. H. (2013). The}
  \ALIUSPlacedTextContent{126.951}{682.096}{382.061}{ALIUSC000000}{\ALIUSFontCormorantRegular}{12.960}{timescale of perceptual evidence integration can be adapted to the}
  \ALIUSPlacedTextContent{126.951}{697.935}{70.808}{ALIUSC000000}{\ALIUSFontCormorantRegular}{12.960}{environment.}
  \ALIUSPlacedTextContent{197.786}{697.935}{88.740}{ALIUSC000000}{\ALIUSFontCormorantItalic}{12.960}{Current Biology. 23}
  \ALIUSPlacedTextContent{286.535}{697.935}{59.448}{ALIUSC000000}{\ALIUSFontCormorantRegular}{12.960}{. 981–986.}
  \ALIUSPlacedTextContent{126.951}{713.536}{198.460}{ALIUSC0000FF}{\ALIUSFontCormorantRegular}{12.960}{https://doi.org/10.1016/j.cub.2013.04.039}
  \ALIUSPlacedTextContent{90.951}{793.143}{122.397}{ALIUSC767171}{\ALIUSFontCormorantLight}{12.000}{ALIUS Bulletin n°4 (2020)}
  \ALIUSPlacedTextContent{387.435}{793.143}{120.728}{ALIUSC767171}{\ALIUSFontCormorantLight}{12.000}{aliusresearch.org/bulletin}
\end{tikzpicture}
\null
\newpage
% Page 13
\thispagestyle{empty}
\noindent\begin{tikzpicture}[remember picture,overlay,x=1bp,y=1bp,shift={(current page.north west)}]
  \fill[ALIUSC0000FF] (155.031bp,-117.351bp) rectangle ++(212.400bp,-0.720bp);
  \fill[ALIUSC0000FF] (297.351bp,-180.231bp) rectangle ++(195.840bp,-0.720bp);
  \fill[ALIUSC0000FF] (126.951bp,-258.951bp) rectangle ++(215.040bp,-0.720bp);
  \fill[ALIUSC0000FF] (126.951bp,-322.071bp) rectangle ++(285.600bp,-0.720bp);
  \fill[ALIUSC0000FF] (126.951bp,-337.671bp) rectangle ++(378.720bp,-0.720bp);
  \fill[ALIUSC0000FF] (126.951bp,-353.511bp) rectangle ++(251.760bp,-0.720bp);
  \fill[ALIUSC0000FF] (126.951bp,-416.391bp) rectangle ++(262.080bp,-0.720bp);
  \fill[ALIUSC0000FF] (126.951bp,-495.111bp) rectangle ++(245.760bp,-0.720bp);
  \fill[ALIUSC0000FF] (126.951bp,-557.991bp) rectangle ++(248.640bp,-0.720bp);
  \ALIUSPlacedTextContent{90.951}{36.329}{321.512}{ALIUSC767171}{\ALIUSFontCormorantRegular}{13.920}{Timothy Hanks – The Neuroscience of Decision Making}
  \ALIUSPlacedTextContent{492.471}{36.329}{16.781}{ALIUSC000000}{\ALIUSFontCormorantRegular}{13.920}{80}
  \ALIUSPlacedTextContent{90.951}{72.976}{417.971}{ALIUSC000000}{\ALIUSFontCormorantRegular}{12.960}{Piet A. T., El Hady A., Brody C. D. (2018). Rats adopt the optimal timescale for}
  \ALIUSPlacedTextContent{126.951}{88.576}{249.218}{ALIUSC000000}{\ALIUSFontCormorantRegular}{12.960}{evidence integration in a dynamic environment.}
  \ALIUSPlacedTextContent{376.988}{88.576}{126.385}{ALIUSC000000}{\ALIUSFontCormorantItalic}{12.960}{Nature Communications. 9}
  \ALIUSPlacedTextContent{503.390}{88.576}{5.593}{ALIUSC000000}{\ALIUSFontCormorantRegular}{12.960}{.}
  \ALIUSPlacedTextContent{126.951}{104.416}{28.045}{ALIUSC000000}{\ALIUSFontCormorantRegular}{12.960}{4265.}
  \ALIUSPlacedTextContent{155.004}{104.416}{212.436}{ALIUSC0000FF}{\ALIUSFontCormorantRegular}{12.960}{https://doi.org/10.1038/s41467-018-06561-y}
  \ALIUSPlacedTextContent{90.951}{135.856}{418.057}{ALIUSC000000}{\ALIUSFontCormorantRegular}{12.960}{Radillo A. E., Veliz-Cuba A., Josić K., Kilpatrick Z. P. (2017). Evidence}
  \ALIUSPlacedTextContent{126.951}{151.696}{381.965}{ALIUSC000000}{\ALIUSFontCormorantRegular}{12.960}{Accumulation and Change Rate Inference in Dynamic Environments.}
  \ALIUSPlacedTextContent{126.951}{167.296}{111.578}{ALIUSC000000}{\ALIUSFontCormorantItalic}{12.960}{Neural Computation. 29}
  \ALIUSPlacedTextContent{238.540}{167.296}{58.768}{ALIUSC000000}{\ALIUSFontCormorantRegular}{12.960}{. 1561–1610.}
  \ALIUSPlacedTextContent{297.312}{167.296}{195.738}{ALIUSC0000FF}{\ALIUSFontCormorantRegular}{12.960}{https://doi.org/10.1162/NECO\_a\_00957}
  \ALIUSPlacedTextContent{90.951}{198.736}{418.033}{ALIUSC000000}{\ALIUSFontCormorantRegular}{12.960}{Scott B. B., Constantinople C. M., Akrami A., Hanks T. D., Brody C. D., Tank D.}
  \ALIUSPlacedTextContent{126.951}{214.576}{382.089}{ALIUSC000000}{\ALIUSFontCormorantRegular}{12.960}{W. (2017) . Fronto-parietal Cortical Circuits Encode Accumulated Evidence}
  \ALIUSPlacedTextContent{126.951}{230.416}{158.010}{ALIUSC000000}{\ALIUSFontCormorantRegular}{12.960}{with a Diversity of Timescales.}
  \ALIUSPlacedTextContent{284.961}{230.416}{50.515}{ALIUSC000000}{\ALIUSFontCormorantItalic}{12.960}{Neuron. 95}
  \ALIUSPlacedTextContent{335.491}{230.416}{65.262}{ALIUSC000000}{\ALIUSFontCormorantRegular}{12.960}{. 385–398.e5.}
  \ALIUSPlacedTextContent{126.951}{246.016}{215.179}{ALIUSC0000FF}{\ALIUSFontCormorantRegular}{12.960}{https://doi.org/10.1016/j.neuron.2017.06.013}
  \ALIUSPlacedTextContent{90.951}{277.456}{418.031}{ALIUSC000000}{\ALIUSFontCormorantRegular}{12.960}{Shadlen M. N, Kiani R, Hanks TD, Churchland A. K. (2008). An intentional}
  \ALIUSPlacedTextContent{126.951}{293.296}{61.029}{ALIUSC000000}{\ALIUSFontCormorantRegular}{12.960}{framework.}
  \ALIUSPlacedTextContent{210.567}{293.296}{34.401}{ALIUSC000000}{\ALIUSFontCormorantRegular}{12.960}{Better}
  \ALIUSPlacedTextContent{267.556}{293.296}{26.083}{ALIUSC000000}{\ALIUSFontCormorantRegular}{12.960}{than}
  \ALIUSPlacedTextContent{316.227}{293.296}{57.126}{ALIUSC000000}{\ALIUSFontCormorantRegular}{12.960}{conscious.;}
  \ALIUSPlacedTextContent{395.942}{293.296}{37.019}{ALIUSC000000}{\ALIUSFontCormorantRegular}{12.960}{71–101.}
  \ALIUSPlacedTextContent{455.538}{293.296}{53.447}{ALIUSC000000}{\ALIUSFontCormorantRegular}{12.960}{Available:}
  \ALIUSPlacedTextContent{126.951}{309.136}{285.537}{ALIUSC0000FF}{\ALIUSFontCormorantRegular}{12.960}{https://books.google.com/books?hl=en\&lr=\&id=JQLg7L-}
  \ALIUSPlacedTextContent{126.951}{324.736}{378.769}{ALIUSC0000FF}{\ALIUSFontCormorantRegular}{12.960}{Qto4C\&oi=fnd\&pg=PA71\&dq=shandlen+decision+making+2008\&ots=Hwh}
  \ALIUSPlacedTextContent{126.951}{340.576}{251.720}{ALIUSC0000FF}{\ALIUSFontCormorantRegular}{12.960}{RIR764A\&sig=BxKVCixK6nd5prL0sh81RSt\_sPM}
  \ALIUSPlacedTextContent{90.951}{372.016}{418.019}{ALIUSC000000}{\ALIUSFontCormorantRegular}{12.960}{Squire R. F, Noudoost B, Schafer R. J, Moore T. (2013). Prefrontal contributions to}
  \ALIUSPlacedTextContent{126.951}{387.856}{144.713}{ALIUSC000000}{\ALIUSFontCormorantRegular}{12.960}{visual selective attention.}
  \ALIUSPlacedTextContent{278.954}{387.856}{171.412}{ALIUSC000000}{\ALIUSFontCormorantItalic}{12.960}{Annual Review Neuroscience. 36}
  \ALIUSPlacedTextContent{450.382}{387.856}{58.599}{ALIUSC000000}{\ALIUSFontCormorantRegular}{12.960}{. 451–466.}
  \ALIUSPlacedTextContent{126.951}{403.456}{262.117}{ALIUSC0000FF}{\ALIUSFontCormorantRegular}{12.960}{https://doi.org/10.1146/annurev-neuro-062111-150439}
  \ALIUSPlacedTextContent{90.951}{434.896}{418.047}{ALIUSC000000}{\ALIUSFontCormorantRegular}{12.960}{Thura D, Cisek P. (2016). Modulation of Premotor and Primary Motor Cortical}
  \ALIUSPlacedTextContent{126.951}{450.736}{382.041}{ALIUSC000000}{\ALIUSFontCormorantRegular}{12.960}{Activity during Volitional Adjustments of Speed-Accuracy Trade-Offs.}
  \ALIUSPlacedTextContent{126.951}{466.336}{108.278}{ALIUSC000000}{\ALIUSFontCormorantItalic}{12.960}{Journal of Neuroscience}
  \ALIUSPlacedTextContent{235.290}{466.336}{5.594}{ALIUSC000000}{\ALIUSFontCormorantRegular}{12.960}{.}
  \ALIUSPlacedTextContent{240.893}{466.336}{9.888}{ALIUSC000000}{\ALIUSFontCormorantItalic}{12.960}{36}
  \ALIUSPlacedTextContent{250.799}{466.336}{59.185}{ALIUSC000000}{\ALIUSFontCormorantRegular}{12.960}{. 938–956.}
  \ALIUSPlacedTextContent{126.951}{482.176}{245.785}{ALIUSC0000FF}{\ALIUSFontCormorantRegular}{12.960}{https://doi.org/10.1523/JNEUROSCI.2230-15.2016}
  \ALIUSPlacedTextContent{90.951}{513.616}{417.985}{ALIUSC000000}{\ALIUSFontCormorantRegular}{12.960}{Wong K-F, Wang X-J. (2006). A recurrent network mechanism of time integration}
  \ALIUSPlacedTextContent{126.951}{529.456}{121.710}{ALIUSC000000}{\ALIUSFontCormorantRegular}{12.960}{in perceptual decisions.}
  \ALIUSPlacedTextContent{248.614}{529.456}{124.104}{ALIUSC000000}{\ALIUSFontCormorantItalic}{12.960}{Journal of Neuroscience. 26}
  \ALIUSPlacedTextContent{372.736}{529.456}{58.398}{ALIUSC000000}{\ALIUSFontCormorantRegular}{12.960}{. 1314–1328.}
  \ALIUSPlacedTextContent{126.951}{545.056}{248.594}{ALIUSC0000FF}{\ALIUSFontCormorantRegular}{12.960}{https://doi.org/10.1523/JNEUROSCI.3733-05.2006}
  \ALIUSPlacedTextContent{90.951}{793.143}{122.397}{ALIUSC767171}{\ALIUSFontCormorantLight}{12.000}{ALIUS Bulletin n°4 (2020)}
  \ALIUSPlacedTextContent{387.435}{793.143}{120.728}{ALIUSC767171}{\ALIUSFontCormorantLight}{12.000}{aliusresearch.org/bulletin}
\end{tikzpicture}
\null
\ifdefined\ALIUSIssueBuild
\clearpage
\expandafter\endinput
\fi
\end{document}
